
\documentclass{article} % For LaTeX2e
\usepackage{iclr2026_conference,times}

% Optional math commands from https://github.com/goodfeli/dlbook_notation.

%% >>> begin math_commands.tex
%%%%% NEW MATH DEFINITIONS %%%%%

\usepackage{amsmath,amsfonts,bm}

% Mark sections of captions for referring to divisions of figures
\newcommand{\figleft}{{\em (Left)}}
\newcommand{\figcenter}{{\em (Center)}}
\newcommand{\figright}{{\em (Right)}}
\newcommand{\figtop}{{\em (Top)}}
\newcommand{\figbottom}{{\em (Bottom)}}
\newcommand{\captiona}{{\em (a)}}
\newcommand{\captionb}{{\em (b)}}
\newcommand{\captionc}{{\em (c)}}
\newcommand{\captiond}{{\em (d)}}

% Highlight a newly defined term
\newcommand{\newterm}[1]{{\bf #1}}


% Figure reference, lower-case.
\def\figref#1{figure~\ref{#1}}
% Figure reference, capital. For start of sentence
\def\Figref#1{Figure~\ref{#1}}
\def\twofigref#1#2{figures \ref{#1} and \ref{#2}}
\def\quadfigref#1#2#3#4{figures \ref{#1}, \ref{#2}, \ref{#3} and \ref{#4}}
% Section reference, lower-case.
\def\secref#1{section~\ref{#1}}
% Section reference, capital.
\def\Secref#1{Section~\ref{#1}}
% Reference to two sections.
\def\twosecrefs#1#2{sections \ref{#1} and \ref{#2}}
% Reference to three sections.
\def\secrefs#1#2#3{sections \ref{#1}, \ref{#2} and \ref{#3}}
% Reference to an equation, lower-case.
\def\eqref#1{equation~\ref{#1}}
% Reference to an equation, upper case
\def\Eqref#1{Equation~\ref{#1}}
% A raw reference to an equation---avoid using if possible
\def\plaineqref#1{\ref{#1}}
% Reference to a chapter, lower-case.
\def\chapref#1{chapter~\ref{#1}}
% Reference to an equation, upper case.
\def\Chapref#1{Chapter~\ref{#1}}
% Reference to a range of chapters
\def\rangechapref#1#2{chapters\ref{#1}--\ref{#2}}
% Reference to an algorithm, lower-case.
\def\algref#1{algorithm~\ref{#1}}
% Reference to an algorithm, upper case.
\def\Algref#1{Algorithm~\ref{#1}}
\def\twoalgref#1#2{algorithms \ref{#1} and \ref{#2}}
\def\Twoalgref#1#2{Algorithms \ref{#1} and \ref{#2}}
% Reference to a part, lower case
\def\partref#1{part~\ref{#1}}
% Reference to a part, upper case
\def\Partref#1{Part~\ref{#1}}
\def\twopartref#1#2{parts \ref{#1} and \ref{#2}}

\def\ceil#1{\lceil #1 \rceil}
\def\floor#1{\lfloor #1 \rfloor}
\def\1{\bm{1}}
\newcommand{\train}{\mathcal{D}}
\newcommand{\valid}{\mathcal{D_{\mathrm{valid}}}}
\newcommand{\test}{\mathcal{D_{\mathrm{test}}}}

\def\eps{{\epsilon}}


% Random variables
\def\reta{{\textnormal{$\eta$}}}
\def\ra{{\textnormal{a}}}
\def\rb{{\textnormal{b}}}
\def\rc{{\textnormal{c}}}
\def\rd{{\textnormal{d}}}
\def\re{{\textnormal{e}}}
\def\rf{{\textnormal{f}}}
\def\rg{{\textnormal{g}}}
\def\rh{{\textnormal{h}}}
\def\ri{{\textnormal{i}}}
\def\rj{{\textnormal{j}}}
\def\rk{{\textnormal{k}}}
\def\rl{{\textnormal{l}}}
% rm is already a command, just don't name any random variables m
\def\rn{{\textnormal{n}}}
\def\ro{{\textnormal{o}}}
\def\rp{{\textnormal{p}}}
\def\rq{{\textnormal{q}}}
\def\rr{{\textnormal{r}}}
\def\rs{{\textnormal{s}}}
\def\rt{{\textnormal{t}}}
\def\ru{{\textnormal{u}}}
\def\rv{{\textnormal{v}}}
\def\rw{{\textnormal{w}}}
\def\rx{{\textnormal{x}}}
\def\ry{{\textnormal{y}}}
\def\rz{{\textnormal{z}}}

% Random vectors
\def\rvepsilon{{\mathbf{\epsilon}}}
\def\rvtheta{{\mathbf{\theta}}}
\def\rva{{\mathbf{a}}}
\def\rvb{{\mathbf{b}}}
\def\rvc{{\mathbf{c}}}
\def\rvd{{\mathbf{d}}}
\def\rve{{\mathbf{e}}}
\def\rvf{{\mathbf{f}}}
\def\rvg{{\mathbf{g}}}
\def\rvh{{\mathbf{h}}}
\def\rvu{{\mathbf{i}}}
\def\rvj{{\mathbf{j}}}
\def\rvk{{\mathbf{k}}}
\def\rvl{{\mathbf{l}}}
\def\rvm{{\mathbf{m}}}
\def\rvn{{\mathbf{n}}}
\def\rvo{{\mathbf{o}}}
\def\rvp{{\mathbf{p}}}
\def\rvq{{\mathbf{q}}}
\def\rvr{{\mathbf{r}}}
\def\rvs{{\mathbf{s}}}
\def\rvt{{\mathbf{t}}}
\def\rvu{{\mathbf{u}}}
\def\rvv{{\mathbf{v}}}
\def\rvw{{\mathbf{w}}}
\def\rvx{{\mathbf{x}}}
\def\rvy{{\mathbf{y}}}
\def\rvz{{\mathbf{z}}}

% Elements of random vectors
\def\erva{{\textnormal{a}}}
\def\ervb{{\textnormal{b}}}
\def\ervc{{\textnormal{c}}}
\def\ervd{{\textnormal{d}}}
\def\erve{{\textnormal{e}}}
\def\ervf{{\textnormal{f}}}
\def\ervg{{\textnormal{g}}}
\def\ervh{{\textnormal{h}}}
\def\ervi{{\textnormal{i}}}
\def\ervj{{\textnormal{j}}}
\def\ervk{{\textnormal{k}}}
\def\ervl{{\textnormal{l}}}
\def\ervm{{\textnormal{m}}}
\def\ervn{{\textnormal{n}}}
\def\ervo{{\textnormal{o}}}
\def\ervp{{\textnormal{p}}}
\def\ervq{{\textnormal{q}}}
\def\ervr{{\textnormal{r}}}
\def\ervs{{\textnormal{s}}}
\def\ervt{{\textnormal{t}}}
\def\ervu{{\textnormal{u}}}
\def\ervv{{\textnormal{v}}}
\def\ervw{{\textnormal{w}}}
\def\ervx{{\textnormal{x}}}
\def\ervy{{\textnormal{y}}}
\def\ervz{{\textnormal{z}}}

% Random matrices
\def\rmA{{\mathbf{A}}}
\def\rmB{{\mathbf{B}}}
\def\rmC{{\mathbf{C}}}
\def\rmD{{\mathbf{D}}}
\def\rmE{{\mathbf{E}}}
\def\rmF{{\mathbf{F}}}
\def\rmG{{\mathbf{G}}}
\def\rmH{{\mathbf{H}}}
\def\rmI{{\mathbf{I}}}
\def\rmJ{{\mathbf{J}}}
\def\rmK{{\mathbf{K}}}
\def\rmL{{\mathbf{L}}}
\def\rmM{{\mathbf{M}}}
\def\rmN{{\mathbf{N}}}
\def\rmO{{\mathbf{O}}}
\def\rmP{{\mathbf{P}}}
\def\rmQ{{\mathbf{Q}}}
\def\rmR{{\mathbf{R}}}
\def\rmS{{\mathbf{S}}}
\def\rmT{{\mathbf{T}}}
\def\rmU{{\mathbf{U}}}
\def\rmV{{\mathbf{V}}}
\def\rmW{{\mathbf{W}}}
\def\rmX{{\mathbf{X}}}
\def\rmY{{\mathbf{Y}}}
\def\rmZ{{\mathbf{Z}}}

% Elements of random matrices
\def\ermA{{\textnormal{A}}}
\def\ermB{{\textnormal{B}}}
\def\ermC{{\textnormal{C}}}
\def\ermD{{\textnormal{D}}}
\def\ermE{{\textnormal{E}}}
\def\ermF{{\textnormal{F}}}
\def\ermG{{\textnormal{G}}}
\def\ermH{{\textnormal{H}}}
\def\ermI{{\textnormal{I}}}
\def\ermJ{{\textnormal{J}}}
\def\ermK{{\textnormal{K}}}
\def\ermL{{\textnormal{L}}}
\def\ermM{{\textnormal{M}}}
\def\ermN{{\textnormal{N}}}
\def\ermO{{\textnormal{O}}}
\def\ermP{{\textnormal{P}}}
\def\ermQ{{\textnormal{Q}}}
\def\ermR{{\textnormal{R}}}
\def\ermS{{\textnormal{S}}}
\def\ermT{{\textnormal{T}}}
\def\ermU{{\textnormal{U}}}
\def\ermV{{\textnormal{V}}}
\def\ermW{{\textnormal{W}}}
\def\ermX{{\textnormal{X}}}
\def\ermY{{\textnormal{Y}}}
\def\ermZ{{\textnormal{Z}}}

% Vectors
\def\vzero{{\bm{0}}}
\def\vone{{\bm{1}}}
\def\vmu{{\bm{\mu}}}
\def\vtheta{{\bm{\theta}}}
\def\va{{\bm{a}}}
\def\vb{{\bm{b}}}
\def\vc{{\bm{c}}}
\def\vd{{\bm{d}}}
\def\ve{{\bm{e}}}
\def\vf{{\bm{f}}}
\def\vg{{\bm{g}}}
\def\vh{{\bm{h}}}
\def\vi{{\bm{i}}}
\def\vj{{\bm{j}}}
\def\vk{{\bm{k}}}
\def\vl{{\bm{l}}}
\def\vm{{\bm{m}}}
\def\vn{{\bm{n}}}
\def\vo{{\bm{o}}}
\def\vp{{\bm{p}}}
\def\vq{{\bm{q}}}
\def\vr{{\bm{r}}}
\def\vs{{\bm{s}}}
\def\vt{{\bm{t}}}
\def\vu{{\bm{u}}}
\def\vv{{\bm{v}}}
\def\vw{{\bm{w}}}
\def\vx{{\bm{x}}}
\def\vy{{\bm{y}}}
\def\vz{{\bm{z}}}

% Elements of vectors
\def\evalpha{{\alpha}}
\def\evbeta{{\beta}}
\def\evepsilon{{\epsilon}}
\def\evlambda{{\lambda}}
\def\evomega{{\omega}}
\def\evmu{{\mu}}
\def\evpsi{{\psi}}
\def\evsigma{{\sigma}}
\def\evtheta{{\theta}}
\def\eva{{a}}
\def\evb{{b}}
\def\evc{{c}}
\def\evd{{d}}
\def\eve{{e}}
\def\evf{{f}}
\def\evg{{g}}
\def\evh{{h}}
\def\evi{{i}}
\def\evj{{j}}
\def\evk{{k}}
\def\evl{{l}}
\def\evm{{m}}
\def\evn{{n}}
\def\evo{{o}}
\def\evp{{p}}
\def\evq{{q}}
\def\evr{{r}}
\def\evs{{s}}
\def\evt{{t}}
\def\evu{{u}}
\def\evv{{v}}
\def\evw{{w}}
\def\evx{{x}}
\def\evy{{y}}
\def\evz{{z}}

% Matrix
\def\mA{{\bm{A}}}
\def\mB{{\bm{B}}}
\def\mC{{\bm{C}}}
\def\mD{{\bm{D}}}
\def\mE{{\bm{E}}}
\def\mF{{\bm{F}}}
\def\mG{{\bm{G}}}
\def\mH{{\bm{H}}}
\def\mI{{\bm{I}}}
\def\mJ{{\bm{J}}}
\def\mK{{\bm{K}}}
\def\mL{{\bm{L}}}
\def\mM{{\bm{M}}}
\def\mN{{\bm{N}}}
\def\mO{{\bm{O}}}
\def\mP{{\bm{P}}}
\def\mQ{{\bm{Q}}}
\def\mR{{\bm{R}}}
\def\mS{{\bm{S}}}
\def\mT{{\bm{T}}}
\def\mU{{\bm{U}}}
\def\mV{{\bm{V}}}
\def\mW{{\bm{W}}}
\def\mX{{\bm{X}}}
\def\mY{{\bm{Y}}}
\def\mZ{{\bm{Z}}}
\def\mBeta{{\bm{\beta}}}
\def\mPhi{{\bm{\Phi}}}
\def\mLambda{{\bm{\Lambda}}}
\def\mSigma{{\bm{\Sigma}}}

% Tensor
\DeclareMathAlphabet{\mathsfit}{\encodingdefault}{\sfdefault}{m}{sl}
\SetMathAlphabet{\mathsfit}{bold}{\encodingdefault}{\sfdefault}{bx}{n}
\newcommand{\tens}[1]{\bm{\mathsfit{#1}}}
\def\tA{{\tens{A}}}
\def\tB{{\tens{B}}}
\def\tC{{\tens{C}}}
\def\tD{{\tens{D}}}
\def\tE{{\tens{E}}}
\def\tF{{\tens{F}}}
\def\tG{{\tens{G}}}
\def\tH{{\tens{H}}}
\def\tI{{\tens{I}}}
\def\tJ{{\tens{J}}}
\def\tK{{\tens{K}}}
\def\tL{{\tens{L}}}
\def\tM{{\tens{M}}}
\def\tN{{\tens{N}}}
\def\tO{{\tens{O}}}
\def\tP{{\tens{P}}}
\def\tQ{{\tens{Q}}}
\def\tR{{\tens{R}}}
\def\tS{{\tens{S}}}
\def\tT{{\tens{T}}}
\def\tU{{\tens{U}}}
\def\tV{{\tens{V}}}
\def\tW{{\tens{W}}}
\def\tX{{\tens{X}}}
\def\tY{{\tens{Y}}}
\def\tZ{{\tens{Z}}}


% Graph
\def\gA{{\mathcal{A}}}
\def\gB{{\mathcal{B}}}
\def\gC{{\mathcal{C}}}
\def\gD{{\mathcal{D}}}
\def\gE{{\mathcal{E}}}
\def\gF{{\mathcal{F}}}
\def\gG{{\mathcal{G}}}
\def\gH{{\mathcal{H}}}
\def\gI{{\mathcal{I}}}
\def\gJ{{\mathcal{J}}}
\def\gK{{\mathcal{K}}}
\def\gL{{\mathcal{L}}}
\def\gM{{\mathcal{M}}}
\def\gN{{\mathcal{N}}}
\def\gO{{\mathcal{O}}}
\def\gP{{\mathcal{P}}}
\def\gQ{{\mathcal{Q}}}
\def\gR{{\mathcal{R}}}
\def\gS{{\mathcal{S}}}
\def\gT{{\mathcal{T}}}
\def\gU{{\mathcal{U}}}
\def\gV{{\mathcal{V}}}
\def\gW{{\mathcal{W}}}
\def\gX{{\mathcal{X}}}
\def\gY{{\mathcal{Y}}}
\def\gZ{{\mathcal{Z}}}

% Sets
\def\sA{{\mathbb{A}}}
\def\sB{{\mathbb{B}}}
\def\sC{{\mathbb{C}}}
\def\sD{{\mathbb{D}}}
% Don't use a set called E, because this would be the same as our symbol
% for expectation.
\def\sF{{\mathbb{F}}}
\def\sG{{\mathbb{G}}}
\def\sH{{\mathbb{H}}}
\def\sI{{\mathbb{I}}}
\def\sJ{{\mathbb{J}}}
\def\sK{{\mathbb{K}}}
\def\sL{{\mathbb{L}}}
\def\sM{{\mathbb{M}}}
\def\sN{{\mathbb{N}}}
\def\sO{{\mathbb{O}}}
\def\sP{{\mathbb{P}}}
\def\sQ{{\mathbb{Q}}}
\def\sR{{\mathbb{R}}}
\def\sS{{\mathbb{S}}}
\def\sT{{\mathbb{T}}}
\def\sU{{\mathbb{U}}}
\def\sV{{\mathbb{V}}}
\def\sW{{\mathbb{W}}}
\def\sX{{\mathbb{X}}}
\def\sY{{\mathbb{Y}}}
\def\sZ{{\mathbb{Z}}}

% Entries of a matrix
\def\emLambda{{\Lambda}}
\def\emA{{A}}
\def\emB{{B}}
\def\emC{{C}}
\def\emD{{D}}
\def\emE{{E}}
\def\emF{{F}}
\def\emG{{G}}
\def\emH{{H}}
\def\emI{{I}}
\def\emJ{{J}}
\def\emK{{K}}
\def\emL{{L}}
\def\emM{{M}}
\def\emN{{N}}
\def\emO{{O}}
\def\emP{{P}}
\def\emQ{{Q}}
\def\emR{{R}}
\def\emS{{S}}
\def\emT{{T}}
\def\emU{{U}}
\def\emV{{V}}
\def\emW{{W}}
\def\emX{{X}}
\def\emY{{Y}}
\def\emZ{{Z}}
\def\emSigma{{\Sigma}}

% entries of a tensor
% Same font as tensor, without \bm wrapper
\newcommand{\etens}[1]{\mathsfit{#1}}
\def\etLambda{{\etens{\Lambda}}}
\def\etA{{\etens{A}}}
\def\etB{{\etens{B}}}
\def\etC{{\etens{C}}}
\def\etD{{\etens{D}}}
\def\etE{{\etens{E}}}
\def\etF{{\etens{F}}}
\def\etG{{\etens{G}}}
\def\etH{{\etens{H}}}
\def\etI{{\etens{I}}}
\def\etJ{{\etens{J}}}
\def\etK{{\etens{K}}}
\def\etL{{\etens{L}}}
\def\etM{{\etens{M}}}
\def\etN{{\etens{N}}}
\def\etO{{\etens{O}}}
\def\etP{{\etens{P}}}
\def\etQ{{\etens{Q}}}
\def\etR{{\etens{R}}}
\def\etS{{\etens{S}}}
\def\etT{{\etens{T}}}
\def\etU{{\etens{U}}}
\def\etV{{\etens{V}}}
\def\etW{{\etens{W}}}
\def\etX{{\etens{X}}}
\def\etY{{\etens{Y}}}
\def\etZ{{\etens{Z}}}

% The true underlying data generating distribution
\newcommand{\pdata}{p_{\rm{data}}}
% The empirical distribution defined by the training set
\newcommand{\ptrain}{\hat{p}_{\rm{data}}}
\newcommand{\Ptrain}{\hat{P}_{\rm{data}}}
% The model distribution
\newcommand{\pmodel}{p_{\rm{model}}}
\newcommand{\Pmodel}{P_{\rm{model}}}
\newcommand{\ptildemodel}{\tilde{p}_{\rm{model}}}
% Stochastic autoencoder distributions
\newcommand{\pencode}{p_{\rm{encoder}}}
\newcommand{\pdecode}{p_{\rm{decoder}}}
\newcommand{\precons}{p_{\rm{reconstruct}}}

\newcommand{\laplace}{\mathrm{Laplace}} % Laplace distribution

\newcommand{\E}{\mathbb{E}}
\newcommand{\Ls}{\mathcal{L}}
\newcommand{\R}{\mathbb{R}}
\newcommand{\emp}{\tilde{p}}
\newcommand{\lr}{\alpha}
\newcommand{\reg}{\lambda}
\newcommand{\rect}{\mathrm{rectifier}}
\newcommand{\softmax}{\mathrm{softmax}}
\newcommand{\sigmoid}{\sigma}
\newcommand{\softplus}{\zeta}
\newcommand{\KL}{D_{\mathrm{KL}}}
\newcommand{\Var}{\mathrm{Var}}
\newcommand{\standarderror}{\mathrm{SE}}
\newcommand{\Cov}{\mathrm{Cov}}
% Wolfram Mathworld says $L^2$ is for function spaces and $\ell^2$ is for vectors
% But then they seem to use $L^2$ for vectors throughout the site, and so does
% wikipedia.
\newcommand{\normlzero}{L^0}
\newcommand{\normlone}{L^1}
\newcommand{\normltwo}{L^2}
\newcommand{\normlp}{L^p}
\newcommand{\normmax}{L^\infty}

\newcommand{\parents}{Pa} % See usage in notation.tex. Chosen to match Daphne's book.

\DeclareMathOperator*{\argmax}{arg\,max}
\DeclareMathOperator*{\argmin}{arg\,min}

\DeclareMathOperator{\sign}{sign}
\DeclareMathOperator{\Tr}{Tr}
\let\ab\allowbreak

%% <<< end math_commands.tex


\usepackage{hyperref}
\usepackage{url}

% added by zengwei
\usepackage{graphicx}
\usepackage{booktabs}
\usepackage{multirow}
\usepackage{caption}
\usepackage{subcaption}
\usepackage{xcolor}
\usepackage{colortbl}
\newcommand{\yao}[1]{\textcolor[rgb]{0,0,1} {#1}}

% \title{Flow2GAN: Hybrid Flow Matching and GAN with Multi-Resolution Network for One-/Two-step High-Fidelity Audio Generation}
\title{Flow2GAN: Hybrid Flow Matching and GAN with Multi-Resolution Network for Few-step High-Fidelity Audio Generation}


% Authors must not appear in the submitted version. They should be hidden
% as long as the \iclrfinalcopy macro remains commented out below.
% Non-anonymous submissions will be rejected without review.

% \author{Zengwei Yao,
% Wei Kang,
% Xiaoyu Yang,
% Fangjun Kuang,
% Liyong Guo,
% Han Zhu,
% \\
% \textbf{Zengrui Jin, 
% Zhaoqing Li, 
% Long Lin, 
% Daniel Povey}  \\
% Xiaomi Corp., Beijing, China \\
% \texttt{dpovey@xiaomi.com} \\
% }

\author{Zengwei Yao, Wei Kang, Han Zhu, Liyong Guo, Lingxuan Ye, Fangjun Kuang, \\
\textbf{Weiji Zhuang, Zhaoqing Li, Zhifeng Han, Long Lin, Daniel Povey}\\
Xiaomi Corp., Beijing, China \\
\texttt{dpovey@xiaomi.com} \\
}

% \author{Antiquus S.~Hippocampus, Natalia Cerebro \& Amelie P. Amygdale \thanks{ Use footnote for providing further information
% about author (webpage, alternative address)---\emph{not} for acknowledging
% funding agencies.  Funding acknowledgements go at the end of the paper.} \\
% Department of Computer Science\\
% Cranberry-Lemon University\\
% Pittsburgh, PA 15213, USA \\
% \texttt{\{hippo,brain,jen\}@cs.cranberry-lemon.edu} \\
% \And
% Ji Q. Ren \& Yevgeny LeNet \\
% Department of Computational Neuroscience \\
% University of the Witwatersrand \\
% Joburg, South Africa \\
% \texttt{\{robot,net\}@wits.ac.za} \\
% \AND
% Coauthor \\
% Affiliation \\
% Address \\
% \texttt{email}
% }

% The \author macro works with any number of authors. There are two commands
% used to separate the names and addresses of multiple authors: \And and \AND.
%
% Using \And between authors leaves it to \LaTeX{} to determine where to break
% the lines. Using \AND forces a linebreak at that point. So, if \LaTeX{}
% puts 3 of 4 authors names on the first line, and the last on the second
% line, try using \AND instead of \And before the third author name.

\newcommand{\fix}{\marginpar{FIX}}
\newcommand{\new}{\marginpar{NEW}}

\iclrfinalcopy % Uncomment for camera-ready version, but NOT for submission.
\begin{document}


\maketitle

\begin{abstract}
% \vspace{-0.5em}

%% >>> begin abstract.tex
Existing dominant methods for audio generation include Generative Adversarial Networks (GANs) and diffusion-based methods like Flow Matching. GANs suffer from slow convergence
during training, while diffusion methods require multi-step inference that introduces considerable computational overhead. In this work, we introduce Flow2GAN, a two-stage framework that combines Flow Matching training for learning generative capabilities with GAN fine-tuning for efficient few-step inference. Specifically, given audio's unique properties, we first improve Flow Matching for audio modeling through: 1) reformulating the objective as endpoint estimation, avoiding velocity estimation difficulties when involving empty regions; 2) applying spectral energy-based loss scaling to emphasize perceptually salient quieter regions. 
Building on these Flow Matching adaptations, we demonstrate that a further stage of lightweight GAN fine-tuning enables us to obtain few-step (e.g., 1/2/4 steps) generators that produce high-quality audio. In addition, we develop a multi-branch network architecture that processes Fourier coefficients at different time-frequency resolutions, which improves the modeling capabilities compared to prior single-resolution designs. 
Experimental results indicate that our Flow2GAN delivers high-fidelity audio generation from Mel-spectrograms or discrete audio tokens, achieving highly favorable quality-efficiency trade-offs compared to existing state-of-the-art GAN-based and Flow Matching-based methods.
Online demo samples are available at \url{https://flow2gan.github.io},
and the source code is released at \url{https://github.com/k2-fsa/Flow2GAN}. 
% Experimental results indicate that our Flow2GAN delivers high-fidelity audio generation from Mel-spectrograms or discrete audio tokens, outperforming existing state-of-the-art GAN-based and Flow Matching-based methods in both quality and efficiency.

% Deep generative models have demonstrated remarkable progress in audio generation, enabling reconstructing high-quality waveforms from compressed representations. Although Generative Adversarial Networks (GANs) have emerged as a dominant approach with one-step inference, they often suffer from slow convergence and potential mode collapse during training. Recent diffusion-based methods like Flow Matching, which offer stable training and powerful generative capabilities, can achieve GAN-competitive performance, but their multi-step inference introduces considerable computational overhead. In this work, we introduce Flow2GAN, a novel framework that combines the strengths of both paradigms, featuring two-stage training and multi-resolution network design, where the first stage employs Flow Matching for learning robust generative capabilities, and the second stage applies GAN fine-tuning for finer-grained generation with efficient one-/two-step inference. Specifically, given audio's unique properties, we first improve Flow Matching for audio modeling through: 1) reformulating the objective as endpoint estimation, avoiding velocity estimation difficulty in empty regions; 2) applying spectral energy-based loss scaling to emphasize perceptually salient quieter regions. Building on these Flow Matching adaptations, we demonstrate that fast GAN fine-tuning enables us to obtain one-step and two-step generators that produce high-quality audio. In addition, we develop a multi-branch network that processes Fourier coefficients at different time-frequency resolutions, showing substantial improvements in modeling capabilities over prior single-resolution designs. Extensive experiments indicate that our Flow2GAN framework delivers high-fidelity audio generation from Mel-spectrograms or discrete acoustic tokens, outperforming existing state-of-the-art GAN-based and Flow Matching-based methods in both quality and efficiency. 


% Deep generative models have demonstrated remarkable progress in audio generation, enabling reconstructing high-quality waveforms from compressed representations. Although Generative Adversarial Networks (GANs) have emerged as a dominant approach with one-step inference, utilizing sophisticated discriminators to capture multi-grained audio details, they often suffer from slow convergence and potential mode collapse during training. Recent diffusion-based methods like Flow Matching, which offer stable training and powerful generative capabilities, can achieve GAN-competitive performance when enhanced with multi-band waveform processing techniques, but their multi-step inference introduces considerable computational overhead. In this work, we introduce Flow2GAN, a two-stage training strategy that combines the strengths of both paradigms, where Flow Matching training in the first stage provides robust generative capabilities, and GAN fine-tuning in the second stage guides the model toward finer-grained generation with efficient one-/two-step inference.  Specifically, considering the unique characteristics of audio signals, we first improve Flow Matching for audio modeling through: 1) reformulating the objective as endpoint estimation, avoiding velocity estimation difficulty in empty regions; 2) applying spectral energy-based loss scaling to emphasize perceptually salient quieter regions. Building on these Flow Matching adaptations, we demonstrate that lightweight GAN fine-tuning enables us to obtain one-step and two-step generators that produce high-quality audio. Additionally, we develop a multi-branch network architecture that processes Fourier spectral coefficients at different time-frequency resolutions, showing substantial improvements in modeling capabilities over prior single-branch design. Extensive experiments show that our Flow2GAN framework delivers high-fidelity audio generation, outperforming existing state-of-the-art GAN-based and Flow Matching-based methods in both quality and efficiency.


%% <<< end abstract.tex

\end{abstract}
% The abstract paragraph should be indented 1/2~inch (3~picas) on both left and
% right-hand margins. Use 10~point type, with a vertical spacing of 11~points.
% The word \textsc{Abstract} must be centered, in small caps, and in point size 12. Two
% line spaces precede the abstract. The abstract must be limited to one
% paragraph.

\vspace{-1em}
\section{Introduction}
\vspace{-0.8em}

%% >>> begin introduction.tex
Audio generation, also known as neural vocoding, aims to reconstruct high-resolution waveforms from compressed representations such as Mel-spectrograms and discrete audio tokens. It serves as a crucial component in various applications, including text-to-speech (TTS)~\citep{neural_tts,vits,zipvoice} and music synthesis~\citep{neural_audio_synth,adversarial_audio_synth,gansynth}, directly impacting the perceptual experience of end users.

Deep generative models have demonstrated remarkable progress in this task. 
Generative Adversarial Networks (GANs)~\citep{gan,wgan} are the dominant approach~\citep{melgan,hifigan,bigvgan,vocos}, leveraging sophisticated discriminators~\citep{melgan, hifigan,mrd,multi-subband-gan} to capture multi-grained audio details and enabling high-fidelity results with efficient one-step inference. However, adversarial optimization in GANs often leads to slow convergence and the risk of mode collapse~\citep{gan_collapse}. Another line of research explores the use of diffusion models~\citep{difussion_2015,ddpm,ddim} for audio generation~\citep{diffwave,Wavegrad,priorgrad,fregrad}, offering stable training and strong generative performance. 
Recent diffusion-based works that integrate multi-band waveform processing~\citep{mbd, rfwave} further boost diffusion on audio data, enabling them to generate high-quality audio.
% Recent diffusion-based works that integrate multi-band waveform processing have achieved GAN-competitive generation quality.
Nevertheless, the strong results of diffusion models come at the cost of multi-step sampling, resulting in high computational demands. 

In this work, we propose Flow2GAN, which follows a two-stage training strategy that combines the strengths of both GAN and diffusion paradigms. In the first stage, we train the model with a Flow Matching objective~\citep{fm1,fm2}, a form of diffusion model~\citep{kingma2023understanding}, to learn robust generative capabilities with stable training. In the second stage, we construct few-step generators from the trained model and utilize GAN fine-tuning for finer-grained generation through adversarial learning. 
To be specific, considering audio's distinctive properties, we improve Flow Matching for audio modeling by: 1) reformulating the training objective as direct endpoint prediction, circumventing the challenges of velocity estimation when involving empty regions; 2) incorporating spectral energy-based loss weighting to emphasize perceptually important low-energy regions. 
Building upon the adapted Flow Matching, the trained model is able to generate clearly improved audio with just 2 sampling steps (Figure~\ref{fig:fm_sample}). This provides a strong foundation for the subsequent GAN fine-tuning stage, which effectively guides the initialized few-step models toward further refined audio quality. Moreover, we design a multi-branch ConvNeXt-based~\citep{convnext} network structure that operates on multi-resolution Fourier coefficients, serving as a powerful backbone with enhanced modeling capabilities compared to the previous single-branch version. 
Our experiments show that Flow2GAN generates high-quality audio across both Mel-spectrogram and audio token conditioning, while providing highly favorable quality-efficiency trade-offs compared to state-of-the-art GAN-based and Flow Matching-based methods.
Ablation studies further validate the effectiveness of each proposed component. 




%% <<< end introduction.tex


\vspace{-1em}
\section{Related work}
\vspace{-0.9em}

%% >>> begin related_work.tex
\noindent\textbf{GAN-based audio generation.} GAN-based models consist of two networks: a generator that produces audio waveforms and a discriminator that classifies whether the generated audio is real or synthetic. The key to these methods is designing sophisticated discriminators to capture audio features from different perspectives, thereby guiding the generator to produce more fine-grained and realistic audio through adversarial learning. 
MelGAN~\citep{melgan} employs a Multi-Scale Discriminator (MSD) that extracts features from waveforms at various downsampled resolutions. HiFi-GAN~\citep{hifigan} extends this approach by incorporating a Multi-Period Discriminator (MPD) that learns complementary structures by analyzing different periodic patterns. Univnet~\citep{mrd} utilizes a Multi-Resolution Discriminator (MRD) that operates in the time-frequency domain across multiple spectral resolutions. BigVGAN~\citep{bigvgan} introduces the Snake activation function to provide periodic inductive bias and incorporates low-pass filters for anti-aliasing. Unlike previous works that rely on transposed convolutions for upsampling, Vocos~\citep{vocos} models spectral coefficients using ConvNeXt model~\citep{convnext} and realizes upsampling to waveforms through the Inverse Short-Time Fourier Transform (ISTFT). Our network structure follows Vocos by using ConvNeXt to operate on spectral coefficients, but we extend it to multiple time-frequency resolutions for a more powerful backbone.

\noindent\textbf{Diffusion/Flow Matching-based audio generation.}
Diffusion-based methods learn transition trajectories from noise distribution to data distribution, excelling at stable training and robust generation. 
DiffWave~\citep{diffwave} and WaveGrad~\citep{Wavegrad} first introduce diffusion models to audio synthesis and get GAN-competitive results. PriorGrad~\citep{priorgrad} uses an adaptive prior derived from data statistics to structure the noise, thereby speeding up the diffusion sampling process. 
PeriodWave~\citep{periodwave} utilizes a multi-period Flow Matching estimator to learn different periodic features of the waveform.
Multi-Band Diffusion (MBD)~\citep{mbd} divides the audio into different frequency bands and processes each band with a distinct diffusion model. RFWave~\citep{rfwave} also adopt this multi-band idea, employing Flow Matching to model each band but processing all bands in parallel with a shared model. These multi-band processing approaches employ frequency equalization to reduce the energy mismatch between the Gaussian distribution and audio data distribution in different frequency bands. While this technique is effective for high-frequency modeling, the equalization process relies on normalization statistics calculated on the dataset~\citep{mbd} or accumulated during training~\citep{rfwave}, which poses a potential generalization risk when applied to distributions that differ significantly from those used to compute the original statistics. 
In this work, we improve Flow Matching for audio modeling by reformulating the objective to endpoint estimation and emphasizing loss in perceptually salient regions, and further combine a GAN fine-tuning stage for fast and finer-grained generation. Note that PriorGrad and RFWave also use energy-based loss scaling~\citep{rfwave}, but they operates at the per-frame level, whereas ours is based on both time and frequency dimensions, which is validated to outperform the per-frame version in our reformulated setting (Table~\ref{tab:ablation_flow}).

\noindent\textbf{Accelerating Flow Matching models in audio generation.} A typical way to accelerate diffusion model inference is through distillation~\citep{progressive_distillation,guided_distillation,consistency_model}, using multi-step outputs to guide fewer-step outputs. WaveFM~\citep{wavefm} introduces a tailored method of consistency distillation~\citep{consistency_model} and can generate audio in 1 step without significant quality degradations. 
Another approach is to integrate GAN objectives with adversarial learning using real data~\citep{fastdiff2, periodwave_turbo, revisiting_distillation}. 
% A recent work in audio generation is PeriodWave-Turbo~\citep{periodwave_turbo}, which introduces GAN fine-tuning to Flow-Matching pre-trained models. Our Flow2GAN shares a similar idea, but we build upon our enhanced Flow Matching specialized for audio modeling, and our 2-step generator achieves competitive results compared to their 4-step model, while offering much faster inference.
A recent work in audio generation is PeriodWave-Turbo~\citep{periodwave_turbo}, which introduces GAN fine-tuning to Flow-Matching pre-trained models. Our Flow2GAN shares a similar idea, but we build upon our enhanced Flow Matching specialized for audio modeling, providing stronger pretrained weights to initialize the GAN generators. This results in significantly superior one-step generator after GAN fine-tuning compared to building upon standard Flow Matching (Table~\ref{tab:ablation_flow}). Notably, PeriodWave-Turbo~\citep{periodwave_turbo} does not explore one-step results. 
Shortcut models~\citep{shortcut} condition on both noise level and step size, enabling better few-step generation than consistency models~\citep{consistency_model}. We also investigate this strategy for few-step audio generation and demonstrate that our approach surpasses it.




%% <<< end related_work.tex


\vspace{-1em}
\section{Preliminary: Flow Matching}
\vspace{-0.8em}

%% >>> begin preliminaries.tex
% \noindent\textbf{}
As a popular formulation of diffusion models~\citep{difussion_2015,ddpm,ddim,rezende2015variational,kingma2023understanding}, Flow Matching~\citep{fm1,fm2} addresses generative modeling by learning the velocity fields that transport the noise distribution $p_{\mathrm{noise}}=\mathcal{N}(0, I)$ to the data distribution $p_{\mathrm{data}}$. Given a noise point $x_0 \sim  p_{\mathrm{noise}}$, a data point $x_1 \sim p_{\mathrm{data}}$, and $t \sim \mathcal{U}[0, 1]$, a straight flow path can be defined by linear interpolation: $x_t = (1-t) x_0 + t x_1$, with corresponding velocity $v_t = x_1 - x_0$.
The same intermediate point $x_t$ can arise from different pairs of $(x_0, x_1)$, each yielding a distinct velocity $v_t$. 
Flow Matching thus trains a network $f_{\theta}(x_t, t)$ to estimate the marginal velocity field~\citep{fm1}, that is, the expected velocity conditioned on $x_t$: $v(x_t, t) = \mathbb{E}_{x_0, x_1} [v_t | x_t]$. 
In practice, this is equivalently optimized in a simpler way by fitting the empirical velocity $v_t$ obtained from randomly sampled pairs $(x_0, x_1)$ during training~\citep{fm1}:
\begin{equation}
\label{eq:fm_loss}
\vspace{-0.5em}
\mathcal{L}_{\mathrm{FM}} = \mathbb{E}_{t,x_0,x_1} \left[\left\| f_{\theta}(x_t, t) - v_t \right\|^2\right].
\vspace{-0.3em}
\end{equation}
In inference, new data can be sampled by solving the ordinary differential equation (ODE) using the learned flow model: $\frac{d x_t}{d t}=f_{\theta}(x_t, t)$.
This process is typically approximated by a numerical solver, such as Euler method, over discrete time steps:  
\begin{equation}
\vspace{-0.7em}
\label{eq:fm_ode}
x_{t_{i+1}} = x_{t_i} + (t_{i+1} - t_i)f_{\theta}(x_{t_i}, t_i).
\vspace{-0.8em}
\end{equation}

% \noindent\textbf{GAN for waveform generation.}
%% <<< end preliminaries.tex


\vspace{-0.8em}
\section{Method}
\vspace{-0.8em}

%% >>> begin method.tex

\begin{figure}
    \centering
    \includegraphics[width=0.8\linewidth]{fig/framework}
    \vspace{-0.6em}
	  \caption{Overall framework of Flow2GAN.}
	\label{fig:framework}
    \vspace{-1.8em}
\end{figure}

As illustrated in Figure~\ref{fig:framework}, Flow2GAN follows a two-stage training strategy: the model is first trained with improved Flow Matching tailored for audio modeling (Section~\ref{ssec:improved_fm}); then few-step generators are built with the trained model and guided by GAN fine-tuning toward fine-grained generation (Section~\ref{ssec:gan_ft}). The model utilizes a multi-branch architecture that processes Fourier coefficients at different time-frequency resolutions (Section~\ref{ssec:multi-branch}).  

\vspace{-0.8em}
\subsection{Improved Flow Matching for audio modeling}
\label{ssec:improved_fm}
\vspace{-0.5em}

\begin{figure}
    \centering
    \includegraphics[width=0.8\linewidth]{fig/fm_sample}
    \vspace{-0.8em}
	  \caption{Example illustration of audio generation. Flow Matching estimates velocity $v_t=x_1-x_0$. The improved Flow Matching in Flow2GAN estimates endpoint $x_1$, yielding cleaner result in silent regions with 2 sampling steps. GAN fine-tuning further enhances the result by filling in details.}
	\label{fig:fm_sample}
    \vspace{-1em}
\end{figure}

% \yao{we focus on one-/two-step flow}

When applying Flow Matching to conditional audio generation, the velocity field is parameterized as $f_{\theta}(x_t, t|\mathbf{c})$, where $\mathbf{c}$ denotes the compressed acoustic representations, such as Mel-spectrograms or discrete audio tokens. 
% We argue that, due to the unique characteristics of audio data (see figure~\ref{}), employing Flow Matching in this setting involves two main challenges:
However, the unique characteristics of audio data present two key challenges for Flow Matching in this context: 1) \textbf{Velocity estimation difficulty.} Audio signals often contain silent segments or frequency bands with zero energy (Figure~\ref{fig:fm_sample} (a)). In these regions, the model must accurately estimate $v_t$ to precisely cancel out the noise in $x_0$ and recover the empty target $x_1$. This increases the learning difficulty as the model is required not only to capture the differences $v_t = x_1 - x_0$ in the active regions of the audio signal, but also to approximate $-x_0$ in these empty regions (Figure~\ref{fig:fm_sample} (c)). 
2) \textbf{Loss-perception mismatch.} The Flow Matching loss (Equation~\ref{eq:fm_loss}) uses mean squared error (MSE) criterion, treating the prediction errors uniformly across all signal regions. However, this does not align well with auditory perception principles, where errors in quieter spectral regions tend to be more perceptually noticeable than equivalent errors in louder regions. 

\vspace{-0.1em}
\noindent\textbf{Reformulation to endpoint estimation.} To circumvent
the difficulty of predicting the velocity $v_t$, we reformulate the Flow Matching problem so that the network predicts the endpoint $x_1$ instead: $\hat{x}_1 =g_{\theta}(x_t, t|\mathbf{c})$. 
Specifically, as $v_t = \frac{x_1 - x_t}{1 - t}$, the loss in Equation~\ref{eq:fm_loss} can be rewritten as: 
\begin{equation}
\label{eq:fm_loss_x1}
\vspace{-0.2em}
\mathcal{L}_{\mathrm{FM}}
= \mathbb{E}_{t,x_0,x_1} \left[\left\| \frac{g_{\theta}(x_t, t|\mathbf{c}) - x_t}{1 - t} - \frac{x_1 - x_t}{1 - t} \right\|^2 \right]
= \mathbb{E}_{t,x_0,x_1} \left[\frac{\left\|g_{\theta}(x_t, t|\mathbf{c}) - x_1 \right\|^2}{(1 - t)^2}\right].
\vspace{-0.1em}
\end{equation}
This reformulation can be interpreted as training the network to reconstruct the clean audio $x_1$ from its noisy version $x_t$ at varying noise levels $(1-t)$. This approach provides a more consistent and stable learning objective tailored to audio data, allowing the network to focus on capturing audio-relevant patterns regardless of whether the region is active or empty. 
% The weighting factor $\frac{1}{(1 - t)^2}$ in Equation~\ref{eq:fm_loss_x1} assigns greater emphasis on instances where $x_t$ is closer to $x_1$. 
Empirically, we found that removing the weighting factor $\frac{1}{(1 - t)^2}$ in Equation~\ref{eq:fm_loss_x1} yields slight improvements at both the Flow-Matching stage (2 steps) and after GAN fine-tuning. The reason might be that without this weighting, the model can better emphasize small $t$ values, which aligns well with our setting.
% Empirically, we found that removing this factor leads to a slight performance degradation in the Flow Matching training stage, but this gap vanishes after GAN fine-tuning (Section~\ref{ssec:gan_ft}) \yao{(even slightly better, no drop at 2 step)}. 
For simplicity, we therefore omit this factor, yielding:
\begin{equation}
\label{eq:fm_loss_x1_no_scale}
\vspace{-0.2em}
\mathcal{L}_{\mathrm{FM}}' 
= \mathbb{E}_{t,x_0,x_1} \left[\left\|g_{\theta}(x_t, t|\mathbf{c}) - x_1 \right\|^2\right].
\vspace{-0.2em}
\end{equation}
%\footnote{In our experiments, we clamp the loss weight $\frac{1}{(1 - t)^2}$ to a maximum of 100.0 for training stability.}.
The sampling process in Equation~\ref{eq:fm_ode} is consequently modified as follows: 
\begin{equation}
\label{eq:fm_ode_x1}
\vspace{-0.3em}
x_{t_{i+1}} = x_{t_i} + (t_{i+1} - t_i)\frac{g_{\theta}(x_{t_i}, t_i|\mathbf{c}) - x_{t_i}}{1 - t_i}.
\vspace{-0.3em}
\end{equation}

\noindent\textbf{Spectral energy-adaptive loss scaling.} To address the loss-perception mismatch, we scale the Flow Matching loss inversely proportional to the energy of the reference spectrogram, thereby encouraging the model to emphasize the perceptually salient quieter regions. Specifically, the prediction error is converted into frequency domain: $\mathcal{S}(g_{\theta}(x_t, t|\mathbf{c}) - x_1)$, where $\mathcal{S}(x)=\mathrm{LinFB}(|\mathrm{STFT}(x)|^2)$ denotes the power STFT followed by a linear filterbank transformation for energy smoothing, and is differentiable. The obtained result is then scaled element-wise inversely by the square root of the reference spectrogram’s energy, $\frac{1}{\sqrt{\mathcal{S}(x_1) + \epsilon}}$, with $\epsilon=1e-7$. Therefore, the loss in Equation~\ref{eq:fm_loss_x1_no_scale} is redefined as:
\begin{equation}
\label{eq:fm_loss_x1_scaling}
\vspace{-0.2em}
\mathcal{L}_{\mathrm{FM}}''
= \mathbb{E}_{t,x_0,x_1} \left[ \sum_{i,j} 
\left(\frac{\mathcal{S}(g_{\theta}(x_t, t|\mathbf{c}) - x_1)}{\sqrt{\mathcal{S}(x_1) + \epsilon}}\right)_{i,j} \right].
% \vspace{-0.2em}
\end{equation}
We clamp the loss scale $\frac{1}{\sqrt{\mathcal{S}(x_1) + \epsilon}}$ between 0.01 and 100 for training stability. See Appendix Table~\ref{tab:model_conf} for detailed configurations of $\mathcal{S}(x)$. 
Note that unlike the per-frame energy-based scaling in previous works~\citep{priorgrad,rfwave}, our scaling strategy is more comprehensive as it also accounts for differences along the frequency dimension.

% We believe the target x_1 might "easier" for the network to learn to output than the velocity x_1 - x_0 because in parts of time-frequency space where x_1 is quiet, the velocity x_1 - x_0 would not necessarily be quiet, yet the process of computing x_1 from (t, x_t and v_t) involves cancellation of the noise in x_0; this would require a very precise estimate of v_t to do accurately as we are predicting a quiet thing from a difference of two loud things.  There is a linear relationship between v_t and x_1, assuming t and x_t are known.  Thus, when we change the target from v_t to x_1 we can view it as being a reformulation of the problem, fundamentally differing only in having a different "t"-dependent weighting of the loss.

% related works: PeriodWave, with discrete wavelet transformation; MBD, EQ

% \vspace{-2em}
\subsection{GAN Fine-tuning for Refinement and Few-step Inference}
\label{ssec:gan_ft}
% \vspace{-0.8em}

% Although the adaptations described in Section~\ref{ssec:improved_fm} improve the audio generation quality of the Flow Matching model with 2 sampling steps (see Figure~\ref{fig:fm_sample} from (d) to (e) and Table~\ref{tab:ablation_flow}), the inherent complexity of audio generation still leaves room for further refinement. 
Although the adaptations described in Section~\ref{ssec:improved_fm} improve the audio generation quality of the Flow Matching model with 2 sampling steps (see Figure~\ref{fig:fm_sample} (d) to (e) and Table~\ref{tab:ablation_flow}), the inherent complexity of audio generation still leaves room for further refinement of details, especially in the high-frequency range. 
Moreover, even though we can obtain higher quality results through multi-step sampling (see Table~\ref{tab:ablation_multi_step} in Appendix), this would require considerable computational overhead. To this end, we employ a GAN fine-tuning strategy for two purposes: enhancing the fine-grained audio details by leveraging the existing carefully designed discriminators~\citep{hifigan, mrd} that capture diverse audio patterns from multiple perspectives, and enabling high-fidelity generation with few-step (even only one-step) inference. 

% We construct the one-step and two-step GAN generators, referred to as $G_{\theta}^1(x_0|\mathbf{c})$ and $G_{\theta}^2(x_0|\mathbf{c})$, by forwarding the trained Flow Matching model through the corresponding number of steps: $G_{\theta}^1(x_0|\mathbf{c})=g_{\theta}(x_0,0|\mathbf{c})$, and $G_{\theta}^2(x_0|\mathbf{c})=g_{\theta}(\frac{x_0 + g_{\theta}(x_0, 0|\mathbf{c})}{2},0.5|\mathbf{c})$. 
We construct the few-step GAN generators, referred to as $G_{\theta}^N(x_0|\mathbf{c})$, by forwarding the trained Flow Matching model through $N$ steps following Equation~\ref{eq:fm_ode_x1}. 
Note that generators with different steps are trained and deployed as separate models rather than a single unified model.
Specifically, for $G_{\theta}^N(x_0|\mathbf{c})$ with $N>1$, gradients are backpropagated through both forward passes, enabling end-to-end optimization of the intermediate output from earlier steps. While this approach can extend to more steps, we focus on one-, two-, and four-step variants for computational efficiency. Following~\citep{vocos}, we adopt the multi-period discriminator (MPD)~\citep{hifigan} and the multi-resolution discriminator (MRD)~\citep{mrd} for adversarial training. The generators are fine-tuned using a combination of HingeGAN adversarial loss~\citep{hingegan}, $L1$ feature matching loss, and multi-scale $L1$ Mel-spectrogram reconstruction loss~\citep{dac}. 

% \yao{Detailed loss scales}.

We found that a fast stage of GAN fine-tuning is sufficient to yield a rapid and substantial improvement in audio quality, while additional fine-tuning can provide further gains (see  Table~\ref{tab:ablation_gan}). Figure~\ref{fig:fm_sample} (e) to (f) shows that the fine-grained details are recovered after GAN fine-tuning. Intuitively, this effect arises because the pre-training stage with the improved Flow Matching objective already equips the model with strong generative modeling capabilities, while subsequent fine-tuning efficiently guides the model toward high-fidelity output by iteratively refining details through adversarial learning with real audio samples. 
As shown in Table~\ref{tab:ablation_flow}, it is worth noting that when using standard Flow Matching for pre-training like~\citep{periodwave_turbo}, the one-step model achieves significantly lower performance after GAN fine-tuning, demonstrating the superiority of our improved Flow Matching objective.
This strategy combines the strengths of both paradigms - improved Flow Matching for robust generative capability learning, and GAN for detail enhancement - ultimately enabling high-fidelity audio generation with fast inference. Notably, compared to pure GAN training, it also requires significantly lower training cost.
 % and mitigates the risk of mode collapsing~\citep{gan,gan_collapse}.

% \vspace{-1em}
\subsection{Multi-resolution network structure}
\label{ssec:multi-branch}
% \vspace{-1em}

\begin{figure}
    \centering
    \includegraphics[width=0.9\linewidth]{fig/network}
    \vspace{-0.5em}
	  \caption{Multi-resolution network structure. Our model processes spectral coefficients from STFT at three different time-frequency resolutions, using larger embedding dimensions for shorter sequences. A ConvNeXt-based encoder processes the conditional compressed representation.}
	\label{fig:network}
    \vspace{-1.7em}
\end{figure}

Inspired by Vocos~\citep{vocos}, our model processes Fourier coefficients rather than waveform samples for saving computational cost and memory usage. While Vocos operates at a single time-frequency resolution, we extend to multi-resolution modeling to better capture audio complexity, providing a more powerful backbone for generative learning.

As shown in Figure~\ref{fig:network}, our model comprises three branches, each processing Fourier coefficients at different resolutions. Specifically, in each branch, the input signal is first transformed via STFT to obtain complex Fourier coefficients, whose real and imaginary components are concatenated along the feature dimension and fed into a ConvNeXT model~\citep{convnext} to produce output complex coefficients. These are then converted back to waveform domain via Inverse STFT (ISTFT). The final output is obtained by summing all branch outputs. Since both STFT and ISTFT operations are differentiable, the entire model is optimized end-to-end. Additionally, following~\citep{periodwave}, we incorporate a ConvNeXt-based condition encoder to extract deeper features from input Mel-spectrograms or codec token embeddings, which serve as shared generation conditions across all branches. Note that for Flow Matching inference, this condition encoder requires only one forward pass, with output features reused across multiple sampling steps. To better balance performance and efficiency, we use larger embedding dimensions for the low-frame-rate branch, and smaller embedding dimensions for the two high-frame-rate branches. Detailed model configurations are provided in Appendix~\ref{tab:model_conf}.  

% \yao{should: add block in Appendix}








%% <<< end method.tex


\vspace{-0.6em}
\section{Experiments}
\vspace{-0.8em}

%% >>> begin experiment.tex
To evaluate the effectiveness of our Flow2GAN, we conduct experiments conditioned on Mel-spectrograms and discrete audio tokens from Encodec~\citep{encodec} (Section~\ref{ssec:comp_sota}). We also perform ablation studies with Mel-spectrogram condition to investigate the effect of each proposed technique (Section~\ref{ssec:ablation}) and compare inference speed with other models (Section~\ref{ssec:speed}). Finally, we evaluate performance when serving as a TTS vocoder (Section~\ref{ssec:comp_tts}). 

\vspace{-0.5em}
\subsection{Experimental Setup} 
\label{ssec:exp_setup}
\vspace{-0.5em}

\noindent\textbf{Datasets and evaluation metrics.} For Mel-spectrogram conditioning, we train models on the LibriTTS dataset~\citep{libritts}, which contains 585 hours of English speech at 24kHz sampling rate. 
We report results on the test set for model comparisons and on the dev set for ablation studies.
For objective evaluation, we use Perceptual Evaluation of Speech Quality (PESQ)~\citep{pesq}, ViSQOL~\citep{visqol}, F1 score for voiced/unvoiced classification (V/UV F1), periodicity error (Periodicity). We also include Fréchet Speech Distance (FSD)~\citep{voicebox}, which measures the similarity between distributions of generated and real samples in a feature space from a Wav2Vec2 encoder~\citep{wav2vec2}. 
For subjective evaluation, we use Similarity MOS (SMOS) between generated and reference audios, with scores from 1 (not at all similar) to 5 (extremely similar), and also 5-scale Mean Opinion Score (MOS) to evaluate naturalness with scores from 1 (completely unnatural) to 5 (completely natural). 
% Following~\citep{bigvgan,vocos}, 
% we also test Similarity MOS (SMOS) between generated and reference audios, with scores from 1 (not at all similar) to 5 (extremely similar). 
We compare against the following state-of-the-art models: BigVGAN~\citep{bigvgan}, Vocos~\citep{vocos}, RFWave~\citep{rfwave}, PeriodWave-Turbo~\citep{periodwave_turbo}, and WaveFM~\citep{wavefm}. In ablation studies, we focus on PESQ and ViSQOL. We use log-Mel spectrograms when comparing to other models and Mel spectrograms in ablation studies, with their results comparison presented in Appendix Table~\ref{tab:comp_mel_logmel}.

For audio token conditioning, following~\citep{mbd,rfwave}, we train on a set of universal audio datasets including: Common Voice 7.0~\citep{commonvoice}, clean speech data from DNS Challenge 4~\citep{dns} (speech), MTG-Jamendo~\citep{mtg} (music), and AudioSet~\citep{audioset} and FSD50K~\citep{fsd50k} (sound events). All audio files are resampled to 24kHz. Following~\citep{rfwave}, we evaluate on a unified test dataset\footnote{https://github.com/bfs18/rfwave}, which contains 900 test samples involving speech, vocals, and sound effect. 
For the universal audios, we report PESQ, ViSQOL and FSD as objective metrics, as well as SMOS and MOS as subjective metric. In this setting,  we compare Flow2GAN against Encodec~\citep{encodec}, MBD~\citep{mbd}, RFWave~\citep{rfwave}, and PeriodWave-Turbo~\citep{periodwave_turbo}. 

% \noindent\textbf{Evaluation Metrics.} For 

To evaluate Flow2GAN as a vocoder in Mel-based TTS systems, we combine it with F5-TTS Base~\citep{f5}, a recent zero-shot TTS model, and compare against various state-of-the-art vocoders. Note that F5-TTS provides two official checkpoints corresponding to librosa and torchaudio mel-spectrogram implementations\footnote{https://github.com/SWivid/F5-TTS}. 
We evaluate zero-shot TTS performance on LibriSpeech-PC~\citep{librispeech-pc} test-clean subset. 
For intelligibility, we report WER between the synthesized speech transcription and the input text. Speaker similarity is measured using cosine similarity between WavLM-based~\citep{wavlm} ECAPA-TDNN~\citep{ecapa} embeddings of the prompt and synthesized speech (denoted as SIM-o~\citep{voicebox}). Naturalness is evaluated using network-based MOS (UTMOS)~\citep{utmos}. 

\noindent\textbf{Implementation details.}
Our Flow2GAN is built with ConvNeXt~\citep{convnext} layers, but we replace the normalization with BiasNorm~\citep{zipformer} and use PreLU activation~\citep{prelu}.
Detailed model configurations are provided in Appendix Table~\ref{tab:model_conf}. 
In GAN fine-tuning stage, following~\citep{dac}, we use multi-scale Mel-spectrogram reconstruction loss with window lengths of $\{32, 64, 128, 256, 512, 1024, 2048\}$, hop lengths set to window length / 4, and Mel bins of $\{5, 10, 20, 40, 80, 160, 320\}$.
For Encodec audio token conditioning, during training we randomly choose a bandwidth from $\{12.0, 6.0, 3.0, 1.5\}$ kbps. At inference, the unified model can support all these bandwidths. We train our models with ScaledAdam optimizer~\citep{zipformer}, which offers faster convergence. Mel-spectrogram conditioned models are trained for 92k Flow Matching iterations and 110k GAN fine-tuning iterations on 2 Nvidia H20 GPUs. Audio token conditioned models are trained for 180k Flow Matching iterations (8 H20 GPUs) and 190k GAN fine-tuning iterations (2 H20 GPUs). For the state-of-the-art models for comparison, we use their released checkpoints and recommended inference configurations.

\vspace{-1em}
\subsection{Comparison to state-of-the-art methods}
\label{ssec:comp_sota}

\noindent\textbf{Mel-spectrogram conditioning.} Table~\ref{tab:comp_mel} compares our Flow2GAN models with other state-of-the-art models with Mel-spectrogram conditioning. 
The one-step Flow2GAN outperforms Vocos~\citep{vocos}, RFWave~\citep{rfwave}, and WaveFM~\citep{wavefm} across most metrics, except PESQ. The two-step and four-step variants outperform all other models except for FSD, SMOS, and MOS. The four-step version achieves the best results in PESQ, ViSQOL, V/UV F1, and Periodicity.
% Except for BigVGAN-v2~\citep{bigvgan} and PeriodWave-Turbo~\citep{periodwave_turbo}, our Flow2GAN models outperform other models across most metrics, with the four-step version achieving the best results in PESQ, V/UV F1, and Periodicity. 
As BigVGAN-v2~\citep{bigvgan} is trained on a larger dataset~\footnote{https://github.com/NVIDIA/BigVGAN}, achieving results close to it demonstrates the effectiveness of our models.
% Note that all Flow2GAN versions are significantly faster than BigVGAN-v2 and PeriodWave-Turbo~\citep{periodwave_turbo} (Table~\ref{tab:ablation_speed}). 


\begin{table}[h]
\centering
\caption{Comparison to state-of-the-art methods on LibriTTS test set with Mel-spectrogram conditioning. $^{*}$BigVGAN-v2 is trained on a large-scale dataset.}
\vspace{-0.5em}
\label{tab:comp_mel}
\setlength{\tabcolsep}{6pt} % Default value: 6pt
% \renewcommand{\arraystretch}{1.2}
\resizebox{1\linewidth}{!}{
\begin{tabular}{l|c|ccccccccc}
\toprule
Model & Params (M) & PESQ$\uparrow$ & ViSQOL$\uparrow$ & V/UV F1$\uparrow$ & Periodicity$\downarrow$ & FSD$\downarrow$ & SMOS$\uparrow$ & MOS$\uparrow$\\
\midrule
BigVGAN & 112.4 & 4.241 & 4.964 & 0.969 & 0.071 & 0.022 & 4.47 $\pm$ 0.10 & 4.33 $\pm$ 0.18\\
 BigVGAN-v2$^{*}$ & 112.4 & 4.379 & 4.971 & 0.978 & 0.055 & \textbf{0.014} & \textbf{4.65 $\pm$ 0.11} & \textbf{4.59 $\pm$ 0.10}\\
Vocos & 13.5 & 3.618 & 4.898 & 0.951 & 0.105 & 0.042 & 4.10 $\pm$ 0.17 & 4.38 $\pm$ 0.16\\
%  BigVSAN & & & & & & & \\
 RFWave (10 steps) & 18.1 & 4.220 & 4.772 & 0.957 & 0.098 & 0.412 & 4.24 $\pm$ 0.16 & 4.29 $\pm$ 0.13 \\
% PeriodWave & & & & & & & \\
PeriodWave-Turbo (4 steps) & 70.2 & 4.434 & 4.965 & 0.958 & 0.096 & 0.020 & 4.20 $\pm$ 0.17 & 4.38 $\pm$ 0.17 \\
 WaveFM (1 step) & 19.5 & 3.540 & 4.894 & 0.943 & 0.124
& 0.098 & 3.72 $\pm$ 0.18 & 3.76 $\pm$ 0.18\\
% \midrule
% Flow2GAN, 1-step \textbf{(ours)} & 78.9 & 4.288 & 4.939 & 0.976 & 0.058 & 0.027 & 4.41 $\pm$ 0.13 \\
%  Flow2GAN, 2-step \textbf{(ours)} & 78.9 & 4.454 & 4.966 & 0.981 & 0.045 & 0.019 & 4.53 $\pm$ 0.12\\
% Flow2GAN, 4-step \textbf{(ours)} & 78.9 & \textbf{4.484} & 4.966 & \textbf{0.984} & \textbf{0.039} & 0.017 &  \\
% log-mel
Flow2GAN, 1-step \textbf{(ours)} & 78.9 & 4.189 & 4.957 & 0.975 & 0.063 & 0.028 & 4.44 $\pm$ 0.14 & 4.39 $\pm$ 0.15 \\
 Flow2GAN, 2-step \textbf{(ours)} & 78.9 & 4.440 & 4.979 & 0.983 & 0.044 & 0.023 & 4.53 $\pm$ 0.13 & 4.56 $\pm$ 0.11 \\
Flow2GAN, 4-step \textbf{(ours)} & 78.9 & \textbf{4.484} & \textbf{4.986} & \textbf{0.985} & \textbf{0.037} & 0.016 & 4.60 $\pm$ 0.14 & 4.58 $\pm$ 0.14\\
\bottomrule
\end{tabular}
}
\vspace{-0.5em}
\end{table}

% \begin{table}[h!]
% \centering
% \caption{ViSQOL$\uparrow$ (audio-mode) comparison on MUSDB18-HQ test set with Mel-spectrogram conditioning. Models are trained on LibriTTS to evaluate out-of-domain performance. $^{*}$BigVGAN-v2 is trained on a large-scale dataset.}
% \label{tab:comp_music}
% \setlength{\tabcolsep}{6pt} % Default value: 6pt
% % \renewcommand{\arraystretch}{1.2}
% \resizebox{0.98\linewidth}{!}{
% \begin{tabular}{l|c|ccccc|c}
% \toprule
% Method & Params (M) & Vocals & Drums & Bass & Other & Mixture & Average \\ 
% \midrule
% BigVGAN~\citep{bigvgan} & & & & & & & \\
% BigVGAN-v2$^{*}$~\citep{bigvgan} & & & & & & & \\
% Vocos~\citep{vocos} & & & & & & & \\
% RFWave~\citep{rfwave} & & & & & & & \\
% PeriodWave~\citep{periodwave} & &  & & & & & \\
% PeriodWave-Turbo~\citep{periodwave_turbo} & & & & & & & \\
% PeriodWave-Turbo~\citep{periodwave_turbo} & & & & & & & \\
% WaveFM~\citep{wavefm} & & & & & & & \\
% \midrule
% Flow2GAN, 1-step \textbf{(ours)} & 78.9 & 3.775 & 3.466 & 4.464 & 3.736 & 3.232 & 3.735 \\
% Flow2GAN, 2-step \textbf{(ours)} & 78.9 & 3.776 & 3.467 & 4.492 & 3.739 & 3.239 & 3.743 \\
% \bottomrule
% \end{tabular}
% }
% \end{table}

% \begin{table}[h!]
% \centering
% \caption{Comparison to BigVGAN-v2 on universal evaluation dataset with Mel-spectrogram conditioning.} 
% \label{tab:comp_codec}
% \setlength{\tabcolsep}{6pt} % Default value: 6pt
% % \renewcommand{\arraystretch}{1.2}
% \resizebox{0.8\linewidth}{!}{
% \begin{tabular}{l|c|ccccccc}
% \toprule
% Method & Params (M) & PESQ$\uparrow$ & ViSQOL$\uparrow$ & FSD $\downarrow$ & MOS$\uparrow$ \\
% \midrule
% BigVGAN-v2 & & & & & \\
% \midrule
% Flow2GAN, 1-step & 78.9 & 1.717 & 3.558 & & \\
% Flow2GAN, 2-step & 78.9 & 1.794 & 3.585 & & \\
% \bottomrule
% \end{tabular}
% }
% \end{table}

\begin{table}[h]
\centering
\caption{Comparison to state-of-the-art methods on universal evaluation dataset with Encodec audio token conditioning.} 
\vspace{-0.5em}
\label{tab:comp_codec}
\setlength{\tabcolsep}{6pt} % Default value: 6pt
% \renewcommand{\arraystretch}{1.2}
\resizebox{1\linewidth}{!}{
\begin{tabular}{c|l|c|cccccccc}
\toprule
Bandwidth (kbps) & Model & Params (M) & PESQ$\uparrow$ & ViSQOL$\uparrow$ & FSD $\downarrow$ & SMOS$\uparrow$ & MOS$\uparrow$ \\
\midrule
1.5 & EnCodec & 56.0 & 1.368 & 3.409 & 1.996 & 1.69 $\pm$ 0.14 & 2.09 $\pm$ 0.27\\
 & MBD & 411.0 & 1.457 & 3.030 & 7.734 & 2.67 $\pm$ 0.18 & 2.83 $\pm$ 0.28\\
& RFWave (10-step) & 18.1 & 1.600 & 3.272 & 2.986 & 2.87 $\pm$ 0.23 & 2.91 $\pm$ 0.23\\
 & PeriodWave-Turbo (4-step) & 31.3 & 1.260 & 3.308 & 4.055 & 1.55 $\pm$ 0.16 & 1.47 $\pm$ 0.19  \\
& Flow2GAN, 1-step \textbf{(ours)} & 78.9 & 1.739 & 3.582 & 1.210 & 2.43 $\pm$ 0.20 & 2.83 $\pm$ 0.21 \\
 & Flow2GAN, 2-step \textbf{(ours)} & 78.9 & 1.803 & 3.609 & 1.152 & 3.04 $\pm$ 0.20 & \textbf{3.86 $\pm$ 0.14} \\
& Flow2GAN, 4-step \textbf{(ours)} & 78.9 & \textbf{1.925} & \textbf{3.662} & \textbf{1.069} & \textbf{3.17 $\pm$ 0.19} & 3.40 $\pm$ 0.18\\
\midrule
 3.0 & EnCodec & 56.0 & 1.745 & 3.867 & 1.380 & 3.11 $\pm$ 0.26 & 3.31 $\pm$ 0.23\\
& MBD & 411.0 & 1.958 & 3.523 & 6.133 & 3.36 $\pm$ 0.21 & 4.00 $\pm$ 0.24\\
 & RFWave (10-step) & 18.1 & 2.124 & 3.784 & 2.622 & 3.41 $\pm$ 0.21 & 3.86 $\pm$ 0.23\\
& PeriodWave-Turbo (4-step) & 31.3 & 2.160 & 4.058 & 1.018 & 3.04 $\pm$ 0.17 & 3.16 $\pm$ 0.23\\
 & Flow2GAN, 1-step \textbf{(ours)} & 78.9 & 2.353 & 4.026 & 0.867 & 3.94 $\pm$ 0.14 & 4.00 $\pm$ 0.19\\
& Flow2GAN, 2-step \textbf{(ours)} & 78.9 & 2.442 & 4.049 & 0.843 & \textbf{4.19 $\pm$ 0.15} & 4.07 $\pm$ 0.24\\
 & Flow2GAN, 4-step \textbf{(ours)} & 78.9 & \textbf{2.550} & \textbf{4.091} & \textbf{0.804} & 4.03 $\pm$ 0.16 & \textbf{4.08 $\pm$ 0.22}\\
\midrule
6.0 & EnCodec & 56.0 & 2.259 & 4.156 & 1.239 & 3.77 $\pm$ 0.19 & 3.47 $\pm$ 0.21\\
 & MBD & 411.0 & 2.094 & 3.690 & 6.177 & 3.15 $\pm$ 0.19 & 4.47 $\pm$ 0.21\\
& RFWave (10-step) & 18.1 & 2.663 & 4.120 & 2.574 & 3.68 $\pm$ 0.22 & 4.00 $\pm$ 0.22\\
 & PeriodWave-Turbo (4-step) & 31.3 & \textbf{3.229} & \textbf{4.424} & 0.712 & 4.00 $\pm$ 0.17 & 4.40 $\pm$ 0.21\\
& Flow2GAN, 1-step \textbf{(ours)} & 78.9 & 2.904 & 4.300 & 0.696 & \textbf{4.46 $\pm$ 0.16} & \textbf{4.42 $\pm$ 0.22}\\
 & Flow2GAN, 2-step \textbf{(ours)} & 78.9 & 2.983 & 4.319 & 0.733 & 4.38 $\pm$ 0.14 & 4.31 $\pm$ 0.17\\
& Flow2GAN, 4-step \textbf{(ours)} & 78.9 & 3.089 & 4.351 & \textbf{0.678} & 4.19 $\pm$ 0.12 & 4.38 $\pm$ 0.13\\
\midrule
 12.0 & EnCodec & 56.0 & 2.776 & 4.347 & 1.204 & 3.76 $\pm$ 0.15 & 4.15 $\pm$ 0.15\\
& MBD & 411.0 & - & - & - & - & - \\
 & RFWave (10-step) & 18.1 & 3.117 & 4.348 & 2.152 & 3.75 $\pm$ 0.18 & 4.13 $\pm$ 0.23\\
& PeriodWave-Turbo (4-step) & 31.3 & \textbf{3.579} & \textbf{4.544} & 0.776 & 4.32 $\pm$ 0.16 & \textbf{4.56 $\pm$ 0.16}\\
 & Flow2GAN, 1-step \textbf{(ours)} & 78.9 & 3.389 & 4.489 & 0.632 & \textbf{4.52 $\pm$ 0.13} & 4.53 $\pm$ 0.23\\
& Flow2GAN, 2-step \textbf{(ours)}& 78.9 & 3.457 & 4.505 & 0.600 & 4.39 $\pm$ 0.14 & 4.44 $\pm$ 0.28 \\
 & Flow2GAN, 4-step \textbf{(ours)} & 78.9 & 3.538 & 4.531 & \textbf{0.557} & 4.22 $\pm$ 0.19 & 4.50 $\pm$ 0.15\\
\bottomrule
\end{tabular}
}
\vspace{-0.5em}
\end{table}

\noindent\textbf{Encodec audio token conditioning.} 
Table~\ref{tab:comp_codec} presents a comparison of Flow2GAN models with other state-of-the-art models conditioned on EnCodec audio tokens. 
At lower bandwidths of 1.5 and 3.0 kbps, all Flow2GAN variants outperform other models across most objective metrics, except that the one-step and two-step models underperform PeriodWave-Turbo~\citep{periodwave_turbo} in ViSQOL at 3.0 kbps. At higher bandwidths of 6.0 and 12.0 kbps, although PeriodWave-Turbo achieves the best results in PESQ and ViSQOL, all Flow2GAN variants outperform other models across all objective metrics, with the four-step Flow2GAN achieving notably better FSD results than PeriodWave-Turbo.
% Both the one-step and two-step Flow2GAN models outperform all other models across most metrics, except for the one-step model at a bandwidth of 1.5 kbps. 
% At bandwidths of 6.0 and 12.0 kbps, the one-step model achieves even higher SMOS than the two-step model. 
With increasing bandwidth, all Flow2GAN models can consistently achieve better SMOS and MOS.

\vspace{-0.5em}
\subsection{Ablation studies}
\label{ssec:ablation}
\vspace{-0.5em}

\noindent\textbf{Ablation of the improved Flow Matching.} 
Table~\ref{tab:ablation_flow} presents ablation results for the Flow Matching stage in Flow2GAN with Mel-spectrogram conditioning. Since we aim for few-step generation after GAN fine-tuning, we focus on the 2-step sampling results for the Flow Matching model. We don't compare one-step Flow Matching results as they tend to converge toward the dataset mean~\citep{shortcut}. 
The results indicate that reformulating the standard Flow Matching objective to endpoint prediction improves PESQ by 0.455 and also yields better results after GAN fine-tuning (from Equation~\ref{eq:fm_loss} to Equation~\ref{eq:fm_loss_x1_no_scale}). Additionally, the proposed spectral energy-adaptive loss scaling (Equation~\ref{eq:fm_loss_x1_scaling}) further improves PESQ and ViSQOL substantially and consistently in both stages. 
% If we set number of filterbank in $\mathcal{S}(x)$ to 1, it degrades to per-frame energy-based scaling, as in ~\citep{rfwave}, leading to clear performance drops in both stages. 
Per-frame energy-based scaling, as in~\citep{priorgrad,rfwave}, leads to clear performance drops in both stages, demonstrating the advantages of our energy-based scaling strategy across time and frequency dimensions.
Maintaining the factor $\frac{1}{(1 - t)^2}$ in Equation~\ref{eq:fm_loss_x1} yields worse results, as it emphasizes larger $t$ values, which is not suitable for our few-step goal. The consistent improvements across both stages confirm that our Flow Matching enhancements contribute to the final performance gains, distinguishing our approach from PeriodWave-Turbo~\citep{periodwave_turbo}, which uses standard Flow Matching for pre-training.


\begin{table}[h!]
\centering
\caption{Ablation study of the improved Flow Matching on LibriTTS dev set with Mel-spectrogram conditioning.}
\vspace{-0.5em}
\label{tab:ablation_flow}
\setlength{\tabcolsep}{3pt} % Default value: 6pt
% \renewcommand{\arraystretch}{1.2}
\resizebox{1\linewidth}{!}{
\begin{tabular}{l|cc|cc|cc}
\toprule
\multirow{3}{*}{Method} & \multicolumn{2}{c|}{Flow Matching trained} & \multicolumn{4}{c}{GAN fine-tuned} \\
 & \multicolumn{2}{c|}{(2 steps)} & \multicolumn{2}{c|}{1-step} & \multicolumn{2}{c}{2-step} \\
 & PESQ$\uparrow$ & ViSQOL$\uparrow$ & PESQ$\uparrow$ & ViSQOL$\uparrow$ & PESQ$\uparrow$ & ViSQOL$\uparrow$ \\
\midrule 
Standard Flow Matching & 2.351 & 3.691 & 3.730 & 4.853 & 4.257 & 4.933 \\
 Predict $x_1$, w/o loss scaling & 2.806 & 3.691 & 4.173 & 4.912 & 4.332 & 4.937 \\
Predict $x_1$, w/ per-frame loss scaling & 3.140 & 3.667 & 4.201 & 4.910 & 4.388 & 4.942 \\
%  Predict $x_1$, w/ $L1(\mathcal{S}(g_{\theta}(x_t,t|\mathbf{c})), \mathcal{S}(x_1))$ & & & & & & \\
 Predict $x_1$, w/ loss scaling, w/ loss factor $\frac{1}{(1 - t)^2}$ & 3.226 & 3.726 & 4.195 & 4.932 & 4.430 & 4.962 \\
Predict $x_1$, w/ loss scaling \textbf{(final)} & \textbf{3.469} & \textbf{3.749} & \textbf{4.303} & \textbf{4.942} & \textbf{4.471} & \textbf{4.969} \\
\bottomrule
\end{tabular}
}
\vspace{-0.5em}
\end{table}

\noindent\textbf{Comparison to pure GAN training.}   
Table~\ref{tab:ablation_gan} demonstrates the effectiveness of our two-stage training paradigm. Pure GAN training converges slowly, whereas Flow2GAN achieves significantly better results with lower training costs by combining Flow Matching training with subsequent GAN fine-tuning. Notably, even a fast GAN fine-tuning with 11k iterations yields high-quality generations for both the one-step and two-step generators, validating the benefits of this training strategy.


\begin{table}[h!]
\centering
\caption{Comparison to pure GAN training on LibriTTS dev set with Mel-spectrogram conditioning.}
\vspace{-0.5em}
\label{tab:ablation_gan}
\setlength{\tabcolsep}{3pt} % Default value: 6pt
% \renewcommand{\arraystretch}{1.2}
\resizebox{0.8\linewidth}{!}{
\begin{tabular}{l|c|c|cc}
\toprule
Method & Training iterations & Training hours & PESQ$\uparrow$ & ViSQOL$\uparrow$ \\
\midrule
Pure GAN & 110k & 26 & 3.587 & 4.823 \\
 & 220k & 53 & 3.770 & 4.838 \\
& 330k & 78 & 3.838 & 4.861 \\
 &  660k & 156 & \textbf{3.919} & \textbf{4.888} \\
\midrule
Flow Matching trained (2 steps) & 92k & 50 & 3.469 & 3.749 \\
\cline{2-5}
 \hspace{2mm} + GAN fine-tuned, 1-step & +11k & +2.6 & 4.195 & 4.887 \\
 & +55k & +13 & 4.286 & 4.933 \\
  & +110k & +26 & \textbf{4.303} & \textbf{4.942} \\
\cline{2-5}
 \hspace{2mm} + GAN fine-tuned, 2-step & +11k & +3.2 & 4.375 & 4.917 \\
  & +55k & +16 & 4.457 & 4.961 \\
 & +110k & +32 & \textbf{4.471} & \textbf{4.969} \\
\bottomrule
\end{tabular}
}
\vspace{-0.5em}
\end{table}


\noindent\textbf{Comparison to Shortcut Models.} 
We also investigate Shortcut models~\citep{shortcut} for few-step audio generation. Specifically, we allocate 75\% of each batch to Flow Matching loss and 25\% to self-consistency loss. See Algorithms 1, 2 in~\citep{shortcut} for training and sampling details. As shown in Table~\ref{tab:comp_shortcut}, Shortcut models improve few-step results over standard Flow Matching in both PESQ and ViSQOL. Our Flow2GAN achieves further performance improvements, demonstrating clear advantages.


 \begin{table}[h!]
\centering
\caption{Comparison to Shortcut models on LibriTTS dev set with Mel-spectrogram conditioning.}
\vspace{-0.5em}
\label{tab:comp_shortcut}
\setlength{\tabcolsep}{6pt} % Default value: 6pt
% \renewcommand{\arraystretch}{1.2}
\resizebox{0.5\linewidth}{!}{
\begin{tabular}{l|c|cc}
\toprule
Method & Steps & PESQ$\uparrow$ & ViSQOL$\uparrow$ \\
\midrule
Standard Flow Matching & 1 & 1.132 & 1.801 \\
 & 2 & 2.351 & 3.691 \\
 & 4 & 3.220 & 4.173 \\
\midrule
Shortcut Models & 1 & 3.181 & 4.426 \\
 & 2 & 3.841 & 4.660 \\
 & 4 & 4.122 & 4.781 \\
\midrule
Flow2GAN & 1 & 4.303 & 4.942 \\
& 2 & 4.471 & \textbf{4.969} \\
& 4 & \textbf{4.498} & 4.967 \\
\bottomrule
\end{tabular}
}
\vspace{-0.5em}
\end{table}

\noindent\textbf{Ablation of the multi-resolution network.} As shown in Table~\ref{tab:ablation_network}, our multi-resolution network structure outperforms the single-resolution baseline~\citep{vocos} with comparable parameters in the final two-step system, demonstrating the benefits of modeling multi-resolution spectral features. 
Removing the condition encoder in Flow2GAN causes a performance drop, indicating that learning higher-level conditional features is helpful.

 % in the one-step model, indicating that learning higher-level conditional features are particularly important when generation capability is limited, as in single-step inference.

 
\begin{table}[h!]
\centering
\caption{Ablation study of model structure on LibriTTS dev set with Mel-spectrogram conditioning.}
\vspace{-0.5em}
\label{tab:ablation_network}
\setlength{\tabcolsep}{6pt} % Default value: 6pt
% \renewcommand{\arraystretch}{1.2}
\resizebox{0.7\linewidth}{!}{
\begin{tabular}{l|c|cc}
\toprule
% \multirow{2}{*}{Method} & \multirow{2}{*}{Params (M)} & \multicolumn{2}{c|}{1-step} & \multicolumn{2}{c}{2-step} \\
% & & PESQ$\uparrow$ & ViSQOL$\uparrow$ & PESQ$\uparrow$ & ViSQOL$\uparrow$ \\
Method (2-step) & Params (M) & PESQ$\uparrow$ & ViSQOL$\uparrow$ \\
\midrule
% Single-resolution, w/o condition encoder & 17.3 & 4.111 & 4.811 \\
 Single-resolution, double layers & 79.9 & 4.442 & 4.920 \\
Multi-resolution, w/o condition encoder & 62.4 & 4.417 & 4.966 \\
 Multi-resolution \textbf{(final)} & 78.9 & \textbf{4.471} & \textbf{4.969} \\
\bottomrule
\end{tabular}
}
\vspace{-0.5em}
\end{table}

\subsection{Inference speed}
\label{ssec:speed}

In table~\ref{tab:ablation_speed}, we compare Flow2GAN with state-of-the-art models under Mel-spectrogram conditioning in terms of inference speed, measured by the speed factor relative to real time (xRT). Experiments are conducted on an Intel Xeon Platinum 8457C CPU and an NVIDIA H20 GPU with batch size of 16 and segment length of 1 second. 
Except for Vocos, 
% both the one-step and two-step 
all versions of Flow2GAN achieve significantly faster inference than other models on both CPU and GPU, while even running faster than real time on CPU. This demonstrates that Flow2GAN excels not only in quality but also in efficiency.

\begin{table}[h!]
\centering
\caption{Model generation speed comparison with batch size of 16 and length of 1 second. BigVGAN-v2$^*$ use the specialized CUDA kernel. }
\vspace{-0.5em}
\label{tab:ablation_speed}
\setlength{\tabcolsep}{6pt} % Default value: 6pt
% \renewcommand{\arraystretch}{1.2}
\resizebox{0.55\linewidth}{!}{
\begin{tabular}{l|c|cc}
\toprule
\multirow{2}{*}{Model} & \multirow{2}{*}{Params (M)} & \multicolumn{2}{c}{xRT$\uparrow$} \\
& & CPU & GPU \\
\midrule
BigVGAN & 112.4 & 0.214 & 69.8 \\
BigVGAN-v2 & 112.4 & 0.214 & 69.8 \\
BigVGAN-v2$^*$ & 112.4 & - & 121.9 \\
Vocos & 13.5 & \textbf{387.57 }& \textbf{6440.80} \\
RFWave (10-step) & 18.1 & 0.37 & 158.8 \\
PeriodWave-Turbo (4-step) & 70.24 & 0.12 & 43.70 \\
WaveFM (1-step) & 19.5 & 0.64 & 226.31 \\
Flow2GAN, 1-step \textbf{(ours)} & 78.9 & 4.85 & 851.67 \\
Flow2GAN, 2-step \textbf{(ours)} & 78.9 & 2.46 & 449.26 \\
Flow2GAN, 4-step \textbf{(ours)} & 78.9 & 1.35 & 228.48 \\
\bottomrule
\end{tabular}
}
\vspace{-0.5em}
\end{table}


\subsection{Zero-shot TTS Performance Comparison}
\label{ssec:comp_tts}


As shown in Table~\ref{tab:comp_tts}, our 4-step Flow2GAN achieves the same SIM-o and slightly better UTMOS than PeriodWave-Turbo~\citep{periodwave_turbo} with significantly faster inference (Table~\ref{tab:ablation_speed}), though with higher WER.
Notably, BigVGAN-v2 performs worse than BigVGAN~\citep{bigvgan}, suggesting potential overfitting to ground-truth Mel-spectrograms. We also find that adding small Gaussian noise $0.2 \times \text{rand}() \times \mathcal{N}(0, 1)$ to the conditioning log-Mel during GAN fine-tuning improves both WER and UTMOS, possibly because it makes the model more robust to the imperfect spectrograms from the TTS diffusion model.

\begin{table}[h!]
\centering
\caption{Zero-shot TTS results on LibriSpeech-PC test-clean with F5-TTS Base models (librosa/torchaudio Mels) and various vocoders. $^{*}$BigVGAN-v2 is trained on a large-scale dataset. $^{\dagger}$ denotes adding noise to log-Mel during GAN fine-tuning.}
\vspace{-0.5em}
\label{tab:comp_tts}
\setlength{\tabcolsep}{6pt} % Default value: 6pt
% \renewcommand{\arraystretch}{1.2}
\resizebox{0.75\linewidth}{!}{
\begin{tabular}{l|c|c|ccc}
\toprule
Model & Params (M) & Mel-type & WER(\%)$\downarrow$ & SIM-o$\uparrow$ & UTMOS$\uparrow$ \\
\midrule
BigVGAN & 112.4 & librosa & \textbf{1.78} & 0.67 & 4.05 \\
BigVGAN-v2$^{*}$ & 112.4 & & 1.82 & 0.66 & 3.70 \\
PeriodWave-Turbo (4 steps) & 70.2 &  & 1.8 & \textbf{0.67} & 4.17 \\ 
\midrule
Vocos & 13.5 & torchaudio &  1.91 & 0.65 & 3.90 \\
RFWave (10 steps) & 18.1 &  & 1.87 & 0.66 & 3.63 \\ 
WaveFM (1 step) & 19.5 &  & 2.01 & 0.65 & 3.15 \\ 
% \midrule
Flow2GAN, 1-step \textbf{(ours)} & 78.9 &  & 1.94 & \textbf{0.67} & 3.75 \\ 
Flow2GAN, 2-step \textbf{(ours)} & 78.9 &  & 1.91 & \textbf{0.67} & 4.11 \\ 
Flow2GAN, 4-step \textbf{(ours)} & 78.9 &  & 1.89 & \textbf{0.67} & 4.18 \\ 
Flow2GAN, 1-step$^{\dagger}$  \textbf{(ours)} & 78.9 &  & 1.88 & \textbf{0.67} & 3.83 \\ 
Flow2GAN, 2-step$^{\dagger}$ \textbf{(ours)} & 78.9 &  & 1.88 & \textbf{0.67} & 4.21 \\ 
Flow2GAN, 4-step$^{\dagger}$ \textbf{(ours)} & 78.9 &  & 1.87 & \textbf{0.67} & \textbf{4.26} \\ 
\bottomrule
\end{tabular}
}
\vspace{-0.5em}
\end{table}

%% <<< end experiment.tex


\vspace{-1em}
\section{Conclusion}
\vspace{-0.8em}

%% >>> begin conclusion.tex


In this work, we propose Flow2GAN, a two-stage training framework. In the first stage, Flow Matching is used to learn the generative capability, while in the second stage, GAN fine-tuning refines the details, ultimately enabling few-step generation. 
Specifically, we improve Flow Matching by reformulating to endpoint estimation to mitigate the difficulties of velocity prediction in empty regions, and by introducing spectral energy–adaptive loss scaling to emphasize perceptually salient quieter regions. 
Based on this, we construct few-step generators from the trained model, and demonstrate that incorporating an additional GAN stage can efficiently enhance the results. 
Our backbone adopts a multi-branch network that processes spectral coefficients at different time–frequency resolutions, providing powerful modeling capabilities.  
Experimental results demonstrate that Flow2GAN achieves efficient high-quality audio generation within few steps (e.g., 1, 2, and 4 steps), delivering highly favorable quality-speed trade-offs in few-step generation compared to existing methods.





%% <<< end conclusion.tex


\vspace{-0.8em}
\section*{Reproducibility statement}
\vspace{-0.8em}

%% >>> begin reproducibility_statement.tex
The datasets used in our experiments are publicly available. Implementation details are provided in Section~\ref{ssec:exp_setup}, and model configurations are described in Appendix Section~\ref{ssec:model_conf}. 
To facilitate the reproduction of our results, the source code and pretrained checkpoints are publicly available at \url{https://github.com/k2-fsa/Flow2GAN}.
% Upon acceptance, we will release the code and pretrained checkpoints for reproduction of our results.
%% <<< end reproducibility_statement.tex


% \section{Submission of conference papers to ICLR 2026}

% ICLR requires electronic submissions, processed by
% \url{https://openreview.net/}. See ICLR's website for more instructions.

% If your paper is ultimately accepted, the statement {\tt
%   {\textbackslash}iclrfinalcopy} should be inserted to adjust the
% format to the camera ready requirements.

% The format for the submissions is a variant of the NeurIPS format.
% Please read carefully the instructions below, and follow them
% faithfully.

% \subsection{Style}

% Papers to be submitted to ICLR 2026 must be prepared according to the
% instructions presented here.

% %% Please note that we have introduced automatic line number generation
% %% into the style file for \LaTeXe. This is to help reviewers
% %% refer to specific lines of the paper when they make their comments. Please do
% %% NOT refer to these line numbers in your paper as they will be removed from the
% %% style file for the final version of accepted papers.

% Authors are required to use the ICLR \LaTeX{} style files obtainable at the
% ICLR website. Please make sure you use the current files and
% not previous versions. Tweaking the style files may be grounds for rejection.

% \subsection{Retrieval of style files}

% The style files for ICLR and other conference information are available online at:
% \begin{center}
%    \url{http://www.iclr.cc/}
% \end{center}
% The file \verb+iclr2026_conference.pdf+ contains these
% instructions and illustrates the
% various formatting requirements your ICLR paper must satisfy.
% Submissions must be made using \LaTeX{} and the style files
% \verb+iclr2026_conference.sty+ and \verb+iclr2026_conference.bst+ (to be used with \LaTeX{}2e). The file
% \verb+iclr2026_conference.tex+ may be used as a ``shell'' for writing your paper. All you
% have to do is replace the author, title, abstract, and text of the paper with
% your own.

% The formatting instructions contained in these style files are summarized in
% sections \ref{gen_inst}, \ref{headings}, and \ref{others} below.

% \section{General formatting instructions}
% \label{gen_inst}

% The text must be confined within a rectangle 5.5~inches (33~picas) wide and
% 9~inches (54~picas) long. The left margin is 1.5~inch (9~picas).
% Use 10~point type with a vertical spacing of 11~points. Times New Roman is the
% preferred typeface throughout. Paragraphs are separated by 1/2~line space,
% with no indentation.

% Paper title is 17~point, in small caps and left-aligned.
% All pages should start at 1~inch (6~picas) from the top of the page.

% Authors' names are
% set in boldface, and each name is placed above its corresponding
% address. The lead author's name is to be listed first, and
% the co-authors' names are set to follow. Authors sharing the
% same address can be on the same line.

% Please pay special attention to the instructions in section \ref{others}
% regarding figures, tables, acknowledgments, and references.


% There will be a strict upper limit of \textbf{9 pages} for the main text of the initial submission, with unlimited additional pages for citations. This limit will be expanded to \textbf{10 pages} for rebuttal/camera ready.

% \section{Headings: first level}
% \label{headings}

% First level headings are in small caps,
% flush left and in point size 12. One line space before the first level
% heading and 1/2~line space after the first level heading.

% \subsection{Headings: second level}

% Second level headings are in small caps,
% flush left and in point size 10. One line space before the second level
% heading and 1/2~line space after the second level heading.

% \subsubsection{Headings: third level}

% Third level headings are in small caps,
% flush left and in point size 10. One line space before the third level
% heading and 1/2~line space after the third level heading.

% \section{Citations, figures, tables, references}
% \label{others}

% These instructions apply to everyone, regardless of the formatter being used.

% \subsection{Citations within the text}

% Citations within the text should be based on the \texttt{natbib} package
% and include the authors' last names and year (with the ``et~al.'' construct
% for more than two authors). When the authors or the publication are
% included in the sentence, the citation should not be in parenthesis using \verb|\citet{}| (as
% in ``See \citet{Hinton06} for more information.''). Otherwise, the citation
% should be in parenthesis using \verb|\citep{}| (as in ``Deep learning shows promise to make progress
% towards AI~\citep{Bengio+chapter2007}.'').

% The corresponding references are to be listed in alphabetical order of
% authors, in the \textsc{References} section. As to the format of the
% references themselves, any style is acceptable as long as it is used
% consistently.

% \subsection{Footnotes}

% Indicate footnotes with a number\footnote{Sample of the first footnote} in the
% text. Place the footnotes at the bottom of the page on which they appear.
% Precede the footnote with a horizontal rule of 2~inches
% (12~picas).\footnote{Sample of the second footnote}

% \subsection{Figures}

% All artwork must be neat, clean, and legible. Lines should be dark
% enough for purposes of reproduction; art work should not be
% hand-drawn. The figure number and caption always appear after the
% figure. Place one line space before the figure caption, and one line
% space after the figure. The figure caption is lower case (except for
% first word and proper nouns); figures are numbered consecutively.

% Make sure the figure caption does not get separated from the figure.
% Leave sufficient space to avoid splitting the figure and figure caption.

% You may use color figures.
% However, it is best for the
% figure captions and the paper body to make sense if the paper is printed
% either in black/white or in color.
% \begin{figure}[h]
% \begin{center}
% %\framebox[4.0in]{$\;$}
% \fbox{\rule[-.5cm]{0cm}{4cm} \rule[-.5cm]{4cm}{0cm}}
% \end{center}
% \caption{Sample figure caption.}
% \end{figure}

% \subsection{Tables}

% All tables must be centered, neat, clean and legible. Do not use hand-drawn
% tables. The table number and title always appear before the table. See
% Table~\ref{sample-table}.

% Place one line space before the table title, one line space after the table
% title, and one line space after the table. The table title must be lower case
% (except for first word and proper nouns); tables are numbered consecutively.

% \begin{table}[t]
% \caption{Sample table title}
% \label{sample-table}
% \begin{center}
% \begin{tabular}{ll}
% \multicolumn{1}{c}{\bf PART}  &\multicolumn{1}{c}{\bf DESCRIPTION}
% \\ \hline \\
% Dendrite         &Input terminal \\
% Axon             &Output terminal \\
% Soma             &Cell body (contains cell nucleus) \\
% \end{tabular}
% \end{center}
% \end{table}

% \section{Default Notation}

% In an attempt to encourage standardized notation, we have included the
% notation file from the textbook, \textit{Deep Learning}
% \cite{goodfellow2016deep} available at
% \url{https://github.com/goodfeli/dlbook_notation/}.  Use of this style
% is not required and can be disabled by commenting out
% \texttt{math\_commands.tex}.


% \centerline{\bf Numbers and Arrays}
% \bgroup
% \def\arraystretch{1.5}
% \begin{tabular}{p{1in}p{3.25in}}
% $\displaystyle a$ & A scalar (integer or real)\\
% $\displaystyle \va$ & A vector\\
% $\displaystyle \mA$ & A matrix\\
% $\displaystyle \tA$ & A tensor\\
% $\displaystyle \mI_n$ & Identity matrix with $n$ rows and $n$ columns\\
% $\displaystyle \mI$ & Identity matrix with dimensionality implied by context\\
% $\displaystyle \ve^{(i)}$ & Standard basis vector $[0,\dots,0,1,0,\dots,0]$ with a 1 at position $i$\\
% $\displaystyle \text{diag}(\va)$ & A square, diagonal matrix with diagonal entries given by $\va$\\
% $\displaystyle \ra$ & A scalar random variable\\
% $\displaystyle \rva$ & A vector-valued random variable\\
% $\displaystyle \rmA$ & A matrix-valued random variable\\
% \end{tabular}
% \egroup
% \vspace{0.25cm}

% \centerline{\bf Sets and Graphs}
% \bgroup
% \def\arraystretch{1.5}

% \begin{tabular}{p{1.25in}p{3.25in}}
% $\displaystyle \sA$ & A set\\
% $\displaystyle \R$ & The set of real numbers \\
% $\displaystyle \{0, 1\}$ & The set containing 0 and 1 \\
% $\displaystyle \{0, 1, \dots, n \}$ & The set of all integers between $0$ and $n$\\
% $\displaystyle [a, b]$ & The real interval including $a$ and $b$\\
% $\displaystyle (a, b]$ & The real interval excluding $a$ but including $b$\\
% $\displaystyle \sA \backslash \sB$ & Set subtraction, i.e., the set containing the elements of $\sA$ that are not in $\sB$\\
% $\displaystyle \gG$ & A graph\\
% $\displaystyle \parents_\gG(\ervx_i)$ & The parents of $\ervx_i$ in $\gG$
% \end{tabular}
% \vspace{0.25cm}


% \centerline{\bf Indexing}
% \bgroup
% \def\arraystretch{1.5}

% \begin{tabular}{p{1.25in}p{3.25in}}
% $\displaystyle \eva_i$ & Element $i$ of vector $\va$, with indexing starting at 1 \\
% $\displaystyle \eva_{-i}$ & All elements of vector $\va$ except for element $i$ \\
% $\displaystyle \emA_{i,j}$ & Element $i, j$ of matrix $\mA$ \\
% $\displaystyle \mA_{i, :}$ & Row $i$ of matrix $\mA$ \\
% $\displaystyle \mA_{:, i}$ & Column $i$ of matrix $\mA$ \\
% $\displaystyle \etA_{i, j, k}$ & Element $(i, j, k)$ of a 3-D tensor $\tA$\\
% $\displaystyle \tA_{:, :, i}$ & 2-D slice of a 3-D tensor\\
% $\displaystyle \erva_i$ & Element $i$ of the random vector $\rva$ \\
% \end{tabular}
% \egroup
% \vspace{0.25cm}


% \centerline{\bf Calculus}
% \bgroup
% \def\arraystretch{1.5}
% \begin{tabular}{p{1.25in}p{3.25in}}
% % NOTE: the [2ex] on the next line adds extra height to that row of the table.
% % Without that command, the fraction on the first line is too tall and collides
% % with the fraction on the second line.
% $\displaystyle\frac{d y} {d x}$ & Derivative of $y$ with respect to $x$\\ [2ex]
% $\displaystyle \frac{\partial y} {\partial x} $ & Partial derivative of $y$ with respect to $x$ \\
% $\displaystyle \nabla_\vx y $ & Gradient of $y$ with respect to $\vx$ \\
% $\displaystyle \nabla_\mX y $ & Matrix derivatives of $y$ with respect to $\mX$ \\
% $\displaystyle \nabla_\tX y $ & Tensor containing derivatives of $y$ with respect to $\tX$ \\
% $\displaystyle \frac{\partial f}{\partial \vx} $ & Jacobian matrix $\mJ \in \R^{m\times n}$ of $f: \R^n \rightarrow \R^m$\\
% $\displaystyle \nabla_\vx^2 f(\vx)\text{ or }\mH( f)(\vx)$ & The Hessian matrix of $f$ at input point $\vx$\\
% $\displaystyle \int f(\vx) d\vx $ & Definite integral over the entire domain of $\vx$ \\
% $\displaystyle \int_\sS f(\vx) d\vx$ & Definite integral with respect to $\vx$ over the set $\sS$ \\
% \end{tabular}
% \egroup
% \vspace{0.25cm}

% \centerline{\bf Probability and Information Theory}
% \bgroup
% \def\arraystretch{1.5}
% \begin{tabular}{p{1.25in}p{3.25in}}
% $\displaystyle P(\ra)$ & A probability distribution over a discrete variable\\
% $\displaystyle p(\ra)$ & A probability distribution over a continuous variable, or over
% a variable whose type has not been specified\\
% $\displaystyle \ra \sim P$ & Random variable $\ra$ has distribution $P$\\% so thing on left of \sim should always be a random variable, with name beginning with \r
% $\displaystyle  \E_{\rx\sim P} [ f(x) ]\text{ or } \E f(x)$ & Expectation of $f(x)$ with respect to $P(\rx)$ \\
% $\displaystyle \Var(f(x)) $ &  Variance of $f(x)$ under $P(\rx)$ \\
% $\displaystyle \Cov(f(x),g(x)) $ & Covariance of $f(x)$ and $g(x)$ under $P(\rx)$\\
% $\displaystyle H(\rx) $ & Shannon entropy of the random variable $\rx$\\
% $\displaystyle \KL ( P \Vert Q ) $ & Kullback-Leibler divergence of P and Q \\
% $\displaystyle \mathcal{N} ( \vx ; \vmu , \mSigma)$ & Gaussian distribution %
% over $\vx$ with mean $\vmu$ and covariance $\mSigma$ \\
% \end{tabular}
% \egroup
% \vspace{0.25cm}

% \centerline{\bf Functions}
% \bgroup
% \def\arraystretch{1.5}
% \begin{tabular}{p{1.25in}p{3.25in}}
% $\displaystyle f: \sA \rightarrow \sB$ & The function $f$ with domain $\sA$ and range $\sB$\\
% $\displaystyle f \circ g $ & Composition of the functions $f$ and $g$ \\
%   $\displaystyle f(\vx ; \vtheta) $ & A function of $\vx$ parametrized by $\vtheta$.
%   (Sometimes we write $f(\vx)$ and omit the argument $\vtheta$ to lighten notation) \\
% $\displaystyle \log x$ & Natural logarithm of $x$ \\
% $\displaystyle \sigma(x)$ & Logistic sigmoid, $\displaystyle \frac{1} {1 + \exp(-x)}$ \\
% $\displaystyle \zeta(x)$ & Softplus, $\log(1 + \exp(x))$ \\
% $\displaystyle || \vx ||_p $ & $\normlp$ norm of $\vx$ \\
% $\displaystyle || \vx || $ & $\normltwo$ norm of $\vx$ \\
% $\displaystyle x^+$ & Positive part of $x$, i.e., $\max(0,x)$\\
% $\displaystyle \1_\mathrm{condition}$ & is 1 if the condition is true, 0 otherwise\\
% \end{tabular}
% \egroup
% \vspace{0.25cm}



% \section{Final instructions}
% Do not change any aspects of the formatting parameters in the style files.
% In particular, do not modify the width or length of the rectangle the text
% should fit into, and do not change font sizes (except perhaps in the
% \textsc{References} section; see below). Please note that pages should be
% numbered.

% \section{Preparing PostScript or PDF files}

% Please prepare PostScript or PDF files with paper size ``US Letter'', and
% not, for example, ``A4''. The -t
% letter option on dvips will produce US Letter files.

% Consider directly generating PDF files using \verb+pdflatex+
% (especially if you are a MiKTeX user).
% PDF figures must be substituted for EPS figures, however.

% Otherwise, please generate your PostScript and PDF files with the following commands:
% \begin{verbatim}
% dvips mypaper.dvi -t letter -Ppdf -G0 -o mypaper.ps
% ps2pdf mypaper.ps mypaper.pdf
% \end{verbatim}

% \subsection{Margins in LaTeX}

% Most of the margin problems come from figures positioned by hand using
% \verb+\special+ or other commands. We suggest using the command
% \verb+\includegraphics+
% from the graphicx package. Always specify the figure width as a multiple of
% the line width as in the example below using .eps graphics
% \begin{verbatim}
%    \usepackage[dvips]{graphicx} ...
%    \includegraphics[width=0.8\linewidth]{myfile.eps}
% \end{verbatim}
% or % Apr 2009 addition
% \begin{verbatim}
%    \usepackage[pdftex]{graphicx} ...
%    \includegraphics[width=0.8\linewidth]{myfile.pdf}
% \end{verbatim}
% for .pdf graphics.
% See section~4.4 in the graphics bundle documentation (\url{http://www.ctan.org/tex-archive/macros/latex/required/graphics/grfguide.ps})

% A number of width problems arise when LaTeX cannot properly hyphenate a
% line. Please give LaTeX hyphenation hints using the \verb+\-+ command.

% \subsubsection*{Author Contributions}
% If you'd like to, you may include  a section for author contributions as is done
% in many journals. This is optional and at the discretion of the authors.

% \subsubsection*{Acknowledgments}
% Use unnumbered third level headings for the acknowledgments. All
% acknowledgments, including those to funding agencies, go at the end of the paper.


\bibliography{iclr2026_conference}
\bibliographystyle{iclr2026_conference}

\newpage
\appendix
\section{Appendix}

%% >>> begin appendix.tex
\subsection{Use of Large Language Models (LLMs)}

In this work, we used LLM as a writing assistance tool to improve phrasing, grammar, and sentence clarity.

\subsection{Multi-step sampling result}

Table~\ref{tab:ablation_multi_step} presents the multi-step sampling results in terms of PESQ and ViSQOL for the Flow Matching stage in Flow2GAN. The results show that increasing the number of sampling steps improves generation quality.

\begin{table}[h!]
\centering
\caption{Multi-step results on LibriTTS dev set for the Flow Matching stage of Flow2GAN.}
\label{tab:ablation_multi_step}
\setlength{\tabcolsep}{6pt} % Default value: 6pt
% \renewcommand{\arraystretch}{1.2}
\resizebox{0.4\linewidth}{!}{
\begin{tabular}{l|cc}
\toprule
Inference steps & PESQ$\uparrow$ & ViSQOL$\uparrow$ \\
\midrule
1 & 2.279 & 2.222 \\
2 & 3.469 & 3.749 \\
4 & 3.788 & 4.156 \\
8 & 4.013 & 4.392 \\
16 & 4.166 & 4.553 \\
\bottomrule
\end{tabular}
}
\end{table}



\subsection{Model configuration}
\label{ssec:model_conf}

% Table~\ref{tab:} 

\begin{table}[h!]
\centering
\caption{Detailed model configurations of Flow2GAN with Mel-spectrogram and Encodec audio token conditioning.}
\label{tab:model_conf}
\setlength{\tabcolsep}{2pt} % Default value: 6pt
% \renewcommand{\arraystretch}{1.2}
\resizebox{1\linewidth}{!}{
\begin{tabular}{l|c|c}
\toprule
Configuration & Mel-spectrogram & Encodec audio token \\
\midrule
Input condition & N-FFT=1024,Hop=256,N-Mel=100 & Latent-dim=128,Hop=320\\
Time embedding dimension & 512 & 512 \\
Branch resolutions (N-FFT,Hop) & (512,256), (256,128), (128,64) & (640,320), (320,160), (160,80) \\
Branch embedding dimensions & 768,512,384 & 768,512,384 \\
Branch Layers & 8,8,8 & 8,8,8 \\
Feedforward hidden factor & 3 & 3 \\
Convolution kernel size & 7 & 7 \\
Condition encoder embedding dimension & 512 & 512 \\
Condition encoder layers & 4 & 4 \\
Loss scaling $\mathcal{S}(x)$, (N-FFT,Hop,N-filters) & (1024,256,256) & (1280,320,320) \\
\bottomrule
\end{tabular}
}
\end{table}

\subsection{Results comparison of Mel and log-Mel conditioning}

Table~\ref{tab:comp_mel_logmel} compares Flow2GAN results with Mel and log-Mel conditioning on the LibriTTS test set.

\begin{table}[h!]
\centering
\caption{Comparison of Mel and log-Mel conditioning on LibriTTS test set.}
\label{tab:comp_mel_logmel}
\setlength{\tabcolsep}{6pt} % Default value: 6pt
% \renewcommand{\arraystretch}{1.2}
\resizebox{0.85\linewidth}{!}{
\begin{tabular}{l|c|cccccc}
\toprule
Model & Condition & PESQ$\uparrow$ & ViSQOL$\uparrow$ & V/UV F1$\uparrow$ & Periodicity$\downarrow$ & FSD$\downarrow$ \\
\midrule
Flow2GAN, 1-step & Mel & \textbf{4.288} & 4.939 & \textbf{0.976} & \textbf{0.058} & \textbf{0.027} \\
Flow2GAN, 1-step & log-Mel & 4.189 & \textbf{4.957} & 0.975 & 0.063 & 0.028 \\
\midrule
 Flow2GAN, 2-step & Mel & \textbf{4.454} & 4.966 & 0.981 & 0.045 & \textbf{0.019} \\
  Flow2GAN, 2-step & log-Mel & 4.440 & \textbf{4.979} & \textbf{0.983} & \textbf{0.044} & 0.023 \\
\midrule
Flow2GAN, 4-step  & Mel & \textbf{4.484} & 4.966 & 0.984 & 0.039 & 0.017 \\
Flow2GAN, 4-step  & log-Mel & \textbf{4.484} & \textbf{4.986} & \textbf{0.985} & \textbf{0.037} & \textbf{0.016} \\
\bottomrule
\end{tabular}
}
\end{table}
%% <<< end appendix.tex

% You may include other additional sections here.


\end{document}
